\documentclass[10pt]{article}
\usepackage[english]{babel}
\usepackage[letterpaper,top=2cm,bottom=2cm,left=2.5cm,right=2.5cm,marginparwidth=1.75cm]{geometry}
\usepackage{amsmath}
\usepackage{dsfont}
\usepackage{amsthm}
\usepackage{amssymb}
\usepackage{graphicx,xcolor,float}
\usepackage{subfigure}
\usepackage{epstopdf}
\usepackage[colorlinks=true, allcolors=blue]{hyperref}
\usepackage{cite}
\usepackage{tikz}
\usepackage{tikz-cd}
\usepackage{wrapfig,lipsum}  
\usepackage{xspace}
\usepackage{aligned-overset}
\usepackage{bbm}
\usepackage{booktabs}
%\usepackage{comment}
\usepackage[utf8]{inputenc}

\usepackage{enumitem}
% \bibliographystyle{alpha}

\usepackage{mathtools}
\mathtoolsset{showonlyrefs}
% this makes equation labels appear only if the equation was cross-referenced at least once

\DeclareMathOperator{\bp}{\mathbf{p}}
\DeclareMathOperator{\bP}{\mathbf{P}}
\DeclareMathOperator{\bq}{\mathbf{q}}
\DeclareMathOperator{\bX}{\mathbf{X}}
\DeclareMathOperator{\TV}{\mathrm{TV}}
\DeclareMathOperator{\Fluc}{\mathrm{Fluc}}
\DeclareMathOperator{\Prob}{\mathrm{Prob}}
\DeclareMathOperator{\tot}{\mathrm{tot}}
\DeclareMathOperator{\Var}{\mathrm{Var}}
\DeclareMathOperator{\rem}{\mathrm{rem}}
\newcommand{\term}[1]{\text{\tt{#1}}\xspace}
\newcommand{\Max}{\term{max}}
\newcommand{\Min}{\term{min}}
\newcommand{\lex}{\term{lex}}
\newcommand{\stable}{\term{stable}}
\newcommand{\Pre}{\term{Pre}_{\term{CL}}}
\newcommand{\CL}{\term{CL}}
\newcommand{\for}{\mathrm{forward}}
\newcommand{\back}{\mathrm{backward}}
\newcommand{\Blocks}{\mathsf{Blocks}}
\newcommand{\new}{\mathrm{new}}
\newcommand{\Leb}{\mathrm{Leb}}
\newcommand{\Unif}{\mathrm{Unif}}
\newcommand{\rmleft}{\mathrm{left}}
\newcommand{\rmright}{\mathrm{right}}
\newcommand{\ileft}{\mathrm{I}_\rmleft}
\newcommand{\iright}{\mathrm{I}_\rmright}
\newcommand{\al}{\mathcal{A}_{\rmleft}}
\newcommand{\ar}{\mathcal{A}_{\rmright}}
\newcommand{\tal}{\widetilde{\mathcal{A}}_{\rmleft}}
\newcommand{\tar}{\widetilde{\mathcal{A}}_{\rmright}}
\newcommand{\lamleft}{\lambda_{\rmleft}}
\newcommand{\lamright}{\lambda_{\rmright}}
\newcommand{\Luce}{\mathrm{Luce}}
\newcommand{\p}{p}

\newcommand{\mscomment}[1]{\textup{\textsf{\color{red}[#1]}}}
\newcommand{\jwcomment}[1]{\textup{\textsf{\color{purple}[#1]}}}
\newcommand{\jscomment}[1]{\textup{\textsf{\color{blue}[#1]}}}





%\newcommand{\be}{\begin{equation}}
%\newcommand{\ee}{\end{equation}}
\newcommand{\bea}{\begin{eqnarray}}
\newcommand{\eea}{\end{eqnarray}}
\newcommand{\<}{\langle}
\renewcommand{\>}{\rangle}
% \newcommand{\E}{{\mathbb E}}
\newcommand{\V}{{\mathbb V}}

\newcommand{\wt}{\widetilde}
\newcommand{\op}{\text{op}}
\newcommand{\wh}{\widehat}
\newcommand{\mcl}{\mathcal}
\newcommand{\Hess}{\mathsf{Hess}}


\def\Ito{It\^{o}}
\def\tY{\tilde{Y}}
\def\Vol{{\rm Vol}}
\def\Info{{\rm I}}
\def\generic{{\rm generic}}
\def\bounded{{\rm bounded}}
\def\eigen{{\rm eigen}}
\def\bdd{{\rm bounded}}
\def\aas{\text{a.a.s.~}}
\def\as{\text{a.s.~}}
\def\ie{\text{i.e.~}}
\def\ae{\text{a.e.~}}
\newcommand\eg{{\text{\eg~}}}
\def\iid{\text{i.i.d.~}}
\def\I{{\rm I}}
\def\II{{\rm II}}
\def\III{{\rm III}}
\def\smaxcut{\mbox{\tiny \rm max-cut}}
\def\sminbisec{\mbox{\tiny \rm min-bis}}
\def\sTV{\mbox{\tiny \rm TV}}
\def\sinit{\mbox{\tiny \rm init}}
\def\seq{\mbox{\tiny \rm eq}}
\def\salg{\mbox{\tiny \rm alg}}
\def\Proj{{\sf P}}
\def\Unif{{\sf Unif}}
\def\BI{{\rm I}}
\def\ov{{\overline v}}
\def\oF{{\overline F}}
\def\eps{{\varepsilon}}
\def\normeq{\cong}
\def\rE{{\rm E}}
\def\id{{\boldsymbol{I}}}
\def\bh{\boldsymbol{h}}
\def\Psp{{\mathfrak{P}}}
\def\diver{{\rm div}}
\def\supp{{\rm supp}}
\def\Cont{\mathscrsfs{C}}
\def\Lp{\mathscrsfs{L}}
\def\cuC{\mathscrsfs{C}}
\def\cuD{\mathscrsfs{D}}
\def\cuP{\mathscrsfs{P}}
\def\cuF{\mathscrsfs{F}}
\def\cuE{\mathscrsfs{E}}
\def\cuK{\mathscrsfs{K}}
\def\cuU{\mathscrsfs{U}}
\def\cuL{\mathscrsfs{L}}
\def\cuV{\mathscrsfs{V}}
\def\oDelta{\overline{\Delta}}
\def\ind{{\mathbbm 1}}
\def\toL{\stackrel{L_1}{\to}}

\def\SF{{\sf SF}}
\def\Law{{\rm Law}}
\def\hx{\hat{x}}
\def\hf{\hat{f}}
\def\hrho{\hat{\rho}}
\def\hxi{\hat{\xi}}
\def\hzeta{\widehat{\zeta}}

\def\cond{{\sf c}}
\def\bcond{\underline{\sf c}}
\def\ro{\overline{{\sf c}}}


\def\btheta{{\boldsymbol{\theta}}}
\def\bTheta{{\boldsymbol{\Theta}}}
\def\bt{{\boldsymbol{\theta}}}
\def\btau{{\boldsymbol{\tau}}}
\def\bsigma{{\boldsymbol{\sigma}}}
\def\obsigma{\overline{\boldsymbol{\sigma}}}
\def\osigma{\overline{\sigma}}
\def\bsm{{\boldsymbol{\sigma}}_{\rm M}}
\def\bxi{{\boldsymbol{\xi}}}
\def\bphi{{\boldsymbol{\phi}}}
\def\bfe{{\boldsymbol{e}}}
\def\bferr{{\boldsymbol{err}}}
\def\bap{\overline{p}}
\def\oR{\overline{R}}



\def\bZ{{\boldsymbol{Z}}}
\def\bW{{\boldsymbol{W}}}
\def\bSigma{{\boldsymbol{\Sigma}}}
\def\bSig{{\boldsymbol{\Sigma}}}
\def\bP{{\boldsymbol{P}}}
\def\bN{{\boldsymbol{N}}}
\def\tbN{\tilde{\boldsymbol{N}}}
\def\bPhi{{\boldsymbol{\Phi}}}
\def\tbPhi{\tilde{\boldsymbol{\Phi}}}
\def\bPp{{\boldsymbol{P}}_{\perp}}

\def\rN{{\rm N}}

\def\hgamma{\widehat{\gamma}}

\def\tell{\tilde{\ell}}
\def\tmu{\tilde{\mu}}
\def\dc{{\Delta{\sf c}}}
\def\bc{{\boldsymbol{c}}}
\def\bC{{\boldsymbol{C}}}
\def\bQ{{\boldsymbol{Q}}}
\def\bM{{\boldsymbol{M}}}
\def\obM{\overline{\boldsymbol{M}}}
\def\bV{{\boldsymbol{V}}}
\def\obV{\overline{\boldsymbol{V}}}

\def\Ude{U^{\delta}}
\def\Xde{X^{\delta}}
\def\Zde{Z^{\delta}}
\def\uDelta{\underline{\Delta}}

\def\sb{{\sf b}}
\def\sB{{\sf B}}
\def\GOE{{\sf GOE}}
\def\hq{{\widehat{q}}}
\def\hQ{{\widehat{Q}}}
\def\hB{{\widehat{B}}}
\def\hbQ{{\boldsymbol{\widehat{Q}}}}
\def\hbu{{\boldsymbol{\widehat{u}}}}
\def\bmu{{\boldsymbol{\mu}}}
\def\bg{{\boldsymbol{g}}}
\def\hbg{{\hat \bg}}
\def\bx{{\boldsymbol{x}}}
\def\hnu{\widehat{\nu}}
\def\bzero{{\mathbf 0}}
\def\bone{{\mathbf 1}}
\def\cF{{\mathcal F}}
\def\cG{{\mathcal G}}
\def\cT{{\mathcal T}}
\def\cC{{\mathcal C}}
\def\cE{{\mathcal E}}
\def\cW{{\mathcal W}}
\def\cX{{\mathcal X}}
\def\cU{{\mathcal U}}
\def\cZ{{\mathcal Z}}
\def\hcC{\widehat{\mathcal C}}
\def\blambda{{\boldsymbol \lambda}}

\def\tz{\tilde{z}}
\def\tbz{\tilde{\boldsymbol z}}
\def\tH{\tilde{H}}
\def\hy{\hat{y}}
\def\cO{{\mathcal O}}
\def\s{\sigma}
\def\Z{{\mathbb Z}}
\def\SS{{\mathbb S}}
\def\RS{\mbox{\tiny\rm RS}}

\def\sS{{\mathscr S}}

\def\opt{\mbox{\tiny\rm opt}}
\def\sUB{\mbox{\tiny\rm UB}}
\def\op{{\rm op}}
\def\sLip{\mbox{\tiny\rm Lip}}
\def\good{\text{good}}
\def\bad{\text{bad}}
\def\relu{\mbox{\tiny\rm ReLU}}
\def\unif{\mbox{\tiny\rm unif}}

\def\vbsigma{\vec \bsig}
\def\vbn{\vec \bn}
\def\vbb{{\vec{\bb}}}
\newcommand{\ubq}{\bar{q}}
\newcommand{\lbq}{\underline{q}}
\def\uu{\underline{u}}

\def\valpha{\vec{\alpha}}
\def\viota{\vec{\iota}}
\def\vk{\vec k}
\def\vp{\vec p}
\def\vv{\vec v}
\def\vw{\vec w}
\def\vx{\vec x}
\def\vq{\vec q}
\def\vz{\vec z}
\def\vA{\vec A}
\def\vB{\vec B}
\def\vC{\vec C}
\def\vD{\vec D}
\def\vV{\vec V}
\def\vX{\vec X}
\def\vW{\vec W}
\def\uvq{\underline{\vec q}}
\def\uvp{\underline{\vec p}}
\def\uvk{\underline{\vec k}}
\def\wtH{\wt{H}}
\def\wtG{\wt{G}}
\def\wtM{\wt{M}}
\def\wtB{\wt{B}}
\def\wtp{\wt{p}}
\def\wtq{\wt{q}}
\def\wtq{\wt{q}}
\def\vone{\vec 1}
\def\vzero{\vec 0}
\def\sig{\sigma}
\def\bsig{{\boldsymbol {\sigma}}}
\def\vbsig{{\vec \bsig}}
\def\ubx{{\underline \bx}}
\def\uby{{\underline \by}}
\def\ubsig{{\underline \bsig}}
\def\ubrho{{\underline \brho}}
\def\ubtau{{\underline \btau}}
\def\ubupsilon{{\underline \bupsilon}}
\def\vh{{\vec h}}
\def\vm{{\vec m}}
\def\vM{{\vec M}}
\def\vu{{\vec u}}
\def\vL{{\vec L}}
\def\vlam{{\vec \lambda}}
\def\va{{\vec a}}
\def\vb{{\vec b}}
\def\ve{{\vec e}}
\def\uvbsig{\underline{\vec \bsig}}
\def\ubsig{\underline{\bsig}}
\def\hbsig{{\widehat \bsig}}
\def\hvbsig{{\boldsymbol{\hat \vbsig}}}
\def\uhbsig{\underline{\boldsymbol{\hat \sigma}}}
\def\uhvbsig{\underline{\boldsymbol{\hat \vbsig}}}
\def\brho{{\boldsymbol \rho}}
\def\vrho{{\vec \rho}}
\def\vbrho{{\vec \brho}}
\def\bupsilon{{\boldsymbol \upsilon}}
\def\bfeta{{\boldsymbol \eta}}
\def\veta{{\vec \eta}}
\def\vbfeta{{\vec \bfeta}}
\def\by{{\boldsymbol y}}
\def\vy{{\vec y}}
\def\vby{{\vec \by}}
\def\vbx{{\vec \bx}}
\def\vg{{\vec g}}
\def\tlambda{{\tilde \lambda}}

\def\naturals{{\mathbb N}}
\def\reals{{\mathbb R}}
\def\R{{\mathbb R}}
\def\complex{{\mathbb C}}
\def\integers{{\mathbb Z}}
\def\Z{{\mathbb Z}}

\def\comp{{\mathsf{comp}}}
\def\Ber{{\mathsf{Ber}}}
\def\poly{{\mathrm{poly}}}

\def\mod{{\mathsf{mod}}}
\def\mix{{\mathsf{mix}}}


\def\sg{s_{\mathfrak G}}
\def\normal{{\sf N}}
\def\cnormal{{\sf CN}}
\def\sT{{\sf T}}
\def\Group{{\mathfrak G}}
\def\sg{s_{{\mathfrak G}}}

\newcommand{\indep}{\perp \!\!\! \perp}


\def\oh{\overline{h}}
\def\bv{{\boldsymbol{v}}}
\def\bz{{\boldsymbol{z}}}
\def\bx{{\boldsymbol{x}}}
\def\ba{{\boldsymbol{a}}}
\def\bb{{\boldsymbol{b}}}
\def\bs{{\boldsymbol{s}}}
\def\tbb{\tilde{\boldsymbol{b}}}
\def\bA{\boldsymbol{A}}
\def\bB{\boldsymbol{B}}
\def\bJ{\boldsymbol{J}}
\def\bH{\boldsymbol{H}}
\def\bT{\boldsymbol{T}}
\def\bm{\boldsymbol{m}}
\def\hbm{\hat{\boldsymbol{m}}}
\def\vbm{\boldsymbol{\vec m}}
\def\vbh{\boldsymbol{\vec h}}
\def\va{\vec{a}}
\def\vb{\vec{b}}
\def\vc{\vec{c}}


\def\obeta{\overline \beta}
\def\ubeta{\underline \beta}

\def\ons{\mathbf{ons}}
\def\ONS{\mathbf{ONS}}

\def\bxz{\boldsymbol{x_0}}
\def\hbx{\boldsymbol{\hat{x}}}

\def\Par{{\sf P}}
\def\de{{\rm d}}
\def\tbX{\tilde{\boldsymbol{X}}}
\def\bX{\boldsymbol{X}}
\def\bY{\boldsymbol{Y}}
\def\bW{\boldsymbol{W}}
\def\prob{{\mathbb P}}
\def\MSE{{\rm MSE}}
\def\Ove{{\rm Overlap}}
\def\<{\langle}
\def\>{\rangle}
\def\Tr{{\sf Tr}}
\def\T{{\sf T}}
\def\Ball{{\sf B}}
\def\range{{\rm range}}
% \def\argmax{{\arg\!\max}}
\def\sign{{\rm sign}}
\def\rank{{\rm rank}}
\def\diag{{\rm diag}}
\def\ed{\stackrel{{\rm d}}{=}}
\def\Poisson{{\rm Poisson}}
\def\cH{{\cal H}}
\def\cM{{\cal M}}
\def\cN{{\cal N}}
\def\usigma{{\underline{\bsigma}}}
\def\uD{\underline{D}}
\def\Pp{{\sf P}^{\perp}}
\def\obv{{\boldsymbol{\overline{v}}}}
\def\di{{\partial i}}
\def\dv{{\partial v}}
\def\conv{{\rm conv}}
\def\cY{{\cal Y}}
\def\cV{{\cal V}}
\def\cL{{\cal L}}

\def\cv{{*}}
\def\by{{\boldsymbol{y}}}
\def\tby{{\widetilde \by}}
\def\bw{{\boldsymbol{w}}}
\def\P{\mathbb{P}}
\def\cE{{\mathcal E}}

\def\txi{\tilde{\xi}}

\def\cov{{\rm Cov}}
\def\crit{{\rm crit}}
\def\SUM{\text{Sum}}
% NOTE TO ANDREA:
% Please algorithm and [noend]{algorithmic} in the beginning of the file, instead the following two below.
%\newcommand{\algorithmicrequire}{\textbf{Input:}}
%\newcommand{\algorithmicensure}{\textbf{Output:}}
%\renewcommand{\algorithmicrequire}{\textbf{Input:}}
%\renewcommand{\algorithmicensure}{\textbf{Output:}}

\def\hp{\hat{p}}
\def\toW{\stackrel{W_2}{\longrightarrow}}
\def\toP{\stackrel{p}{\longrightarrow}}
\def\toLtwo{\stackrel{L_2}{\longrightarrow}}
\def\eqP{\stackrel{p}{=}}
\def\tol {{\tt tol}}
\def\maxiter{{\tt maxiter}}
\def\be{{\boldsymbol{e}}}
\def\blambda{{\boldsymbol{\lambda}}}
\def\bD{{\boldsymbol{D}}}
\def\deltaMax{\delta_\text{max}}
\def\DeltaMax{\Delta_\text{max}}
\def\mRS{m_\text{\tiny RSB}}

\def\cD{{\cal D}}
\def\hcD{\widehat{\cal D}}

\def\vE{\vec{E}}
\def\bphi{{\boldsymbol{\varphi}}}
\def\obphi{\overline{\boldsymbol{\varphi}}}
\def\bu{{\boldsymbol{u}}}
\def\b0{{\boldsymbol{0}}}

\def\hG{\widehat{G}}
\def\hH{\widehat{H}}

\def\bhSigma{{\bf\widehat{\Sigma}}} 
\def\Bind{{\sf Bind}}
\def\Bin{{\sf Bin}}
\def\Hyp{{\sf Hyp}}
\def\Mult{{\sf Mult}}
\def\Var{{\rm Var}}
\def\Input{{\rm Input}}
\def\Past{{\rm Past}}

\def\enter{\mbox{\tiny\rm enter}}
\def\exit{\mbox{\tiny\rm exit}}
\def\return{\mbox{\tiny\rm return}}
\def\sBL{\mbox{\tiny\rm BL}}
\def\rad{\overline{\rho}}
\def\rd{{\mathrm {rad}}}
% \def\span{{\mathrm {span}}}
\def\ov{\overline{v}}
\def\ou{\overline{u}}
\def\bfone{{\boldsymbol 1}}
\def\bff{{\boldsymbol f}}
\def\bF{{\boldsymbol F}}
\def\bG{{\boldsymbol G}}
\def\bGG{{\boldsymbol{Div}}}
\def\bR{{\boldsymbol R}}

\def\obw{\overline{\boldsymbol w}}
\def\hbw{\hat{\boldsymbol w}}
\def\onu{\overline{\nu}}
\def\orh{\overline{\rho}}

\def\obtheta{\overline{\boldsymbol \theta}}
\def\tbtheta{\tilde{\boldsymbol \theta}}
\def\tbx{\tilde{\boldsymbol x}}

\def\rP{{\rm P}}

\DeclareMathOperator*{\plim}{p-lim}
\DeclareMathOperator*{\plimsup}{p-limsup}
\DeclareMathOperator*{\pliminf}{p-liminf}

\def\OPT{{\sf OPT}}
\def\hF{\widehat{F}}
\def\tcF{\tilde{\mathcal F}}
\def\oW{\overline{W}}
\def\oq{\overline{q}}
\def\uq{\underline{q}}
\def\oQ{\overline{Q}}
\def\obX{\overline{\boldsymbol X}}
\def\bdelta{{\boldsymbol \delta}}
\def\ocD{{\overline {\mathscrsfs{P}}}}
\def\cD{{{\mathcal D}}}
\def\tcL{{\tilde {\mathcal L}}}
\def\ent{{\rm Ent}}

\def\hg{\widehat{g}}
\def\trho{\tilde{\rho}}

\def\esssup{{\rm esssup}}
\def\essinf{{\rm essinf}}

% \def\root{\o}
\def\tu{\tilde{u}}
\def\tbu{\tilde{\boldsymbol u}}
\def\tU{\tilde{U}}
\def\tV{\tilde{V}}
\def\cA{{\mathcal A}}
\def\cI{{\mathcal I}}
\def\cS{{\mathcal S}}
\def\bn{{\boldsymbol n}}
\def\cB{{\mathcal B}}
\def\br{{\boldsymbol r}}
% \def\argmin{{\rm arg\,min}}
\def\db{{\bar d}}
\def\cK{{\mathcal K}}
\def\ess{{\rm ess}}
% \renewcommand{\l}{\vert}
\def\bgamma{{\boldsymbol \gamma}}
\def\balpha{{\boldsymbol \alpha}}
\def\setE{{E}}
\def\setS{{S}}

\def\err{{\sf err}}
\def\lambdaknot{\lambda_0}
\def\barR{\overline R}
\def\star{*}
\def\rhon{\rho^\infty}
\def\bzero{\boldsymbol{0}}
\def\ozeta{\overline{\zeta}}
\def\olambda{\overline{\lambda}}
\def\uzeta{\underline{\zeta}}
\def\oI{\overline{I}}
\def\uI{\underline{I}}
\def\oJ{\overline{J}}
\def\uJ{\underline{J}}
\def\uK{\underline{K}}
\def\uQ{\underline{Q}}
\def\ul{\underline{\ell}}
\def\ol{\overline{\ell}}
\def\PA{\mathrm{pa}}

\newcommand{\x}{\mathbf{x}}
\newcommand{\y}{\mathbf{y}}
\newcommand{\z}{\mathbf{z}}
\newcommand{\w}{\mathbf{w}}
\newcommand{\f}{\mathbf{f}}
% \renewcommand{\u}{\mathbf{u}}
% \renewcommand{\v}{\mathbf{v}}
\renewcommand{\b}{\mathbf{b}}
\newcommand{\m}{\mathbf{m}}

%\newcommand{\val}{\texttt{val}}

\def\beps{{\boldsymbol \eps}}

\def\fr{\frac}
\def\lt{\left}
\def\rt{\right}

\def\la{\langle}
\def\ra{\rangle}

\def\eps{\varepsilon}

\def\bbA{{\mathbb{A}}}
\def\bbB{{\mathbb{B}}}
\def\bbC{{\mathbb{C}}}
\def\bbE{{\mathbb{E}}}
\def\bbH{{\mathbb{H}}}
\def\bbI{{\mathbb{I}}}
\def\hbbI{{\hat{\mathbb{I}}}}
\def\bbL{{\mathbb{L}}}
\def\bbN{{\mathbb{N}}}
\def\bbP{{\mathbb{P}}}
\def\bbR{{\mathbb{R}}}
\def\bbS{{\mathbb{S}}}
\def\bbT{{\mathbb{T}}}
\def\bbW{{\mathbb{W}}}
\def\bbZ{{\mathbb{Z}}}

\def\tbbI{{\widetilde \bbI}}

\def\ubbL{\underline{\mathbb{L}}}

\def\cA{{\mathcal{A}}}
\def\cB{{\mathcal{B}}}
\def\cF{{\mathcal{F}}}
\def\cK{{\mathcal{K}}}
\def\cN{{\mathcal{N}}}
\def\cP{{\mathcal{P}}}
\def\cQ{{\mathcal{Q}}}
\def\cR{{\mathcal{R}}}
\def\ucQ{\underline{\mathcal{Q}}}

\def\Crt{{\mathsf{Crt}}}
\def\tot{{\text{tot}}}
\def\unstable{{\mathsf{unstable}}}

\def\sC{{\mathscr{C}}}
\def\sH{{\mathscr{H}}}
\def\sK{{\mathscr{K}}}
\def\sN{{\mathscr{N}}}
\def\sM{{\mathscr{M}}}
\def\sU{{\mathscr{U}}}

\def\one{{\mathbbm{1}}}
\def\bg{{\mathbf{g}}}
\def\bh{{\boldsymbol{h}}}
\def\bk{{\mathbf{k}}}
\def\bp{{\mathbf{p}}}
\def\bq{{\mathbf{q}}}
\def\hbsig{{\boldsymbol {\hat \sigma}}}
\def\bxi{{\boldsymbol \xi}}
\def\bHN{\mathbf{H}_N}
\def\bM{\mathbf{M}}
\def\bP{\mathbf{P}}
\def\bQ{\mathbf{Q}}
\def\bR{\mathbf{R}}
\def\bS{\mathbf{S}}
\def\hcA{{\widehat \cA}}
\def\hcP{{\hat \cP}}
\def\hcK{{\hat \cK}}

\def\alg{{\mathsf{alg}}}
\def\ALG{{\mathsf{ALG}}}
\def\hALG{{\widehat{\mathsf{ALG}}}}
\def\oALG{\overline{\mathsf{ALG}}}
\def\OPT{{\mathsf{OPT}}}
\def\OGP{{\mathsf{OGP}}}
\def\BOGP{{\mathsf{BOGP}}}

\def\TV{{\mathrm{TV}}}
\def\GS{{\mathrm{GS}}}
\def\bGS{{\boldsymbol{\mathrm{GS}}}}

\def\Ssolve{S_{\mathrm{solve}}}
\def\Sconc{S_{\mathrm{conc}}}
\def\Soverlap{S_{\mathrm{overlap}}}
\def\Soffset{S_{\mathrm{offset}}}
\def\Sogp{S_{\mathrm{ogp}}}
\def\Seigen{S_{\mathrm{eigen}}}
\def\vsig{\vec \sigma}
\def\pmax{{p_{\max}}}
\def\smax{{s_{\max}}}
\def\lammax{{\lambda_{\max}}}

\def\psist{{\psi^*}}

\DeclareMathOperator*{\E}{\bbE}
\DeclareMathOperator*{\argmax}{\arg\max}
\DeclareMathOperator*{\argmin}{\arg\min}
% \def\argmax{{\rm arg\,max}}

\def\Sp{\mathrm{Sp}}
\def\Is{\mathrm{Is}}

\def\sph{\mathrm{sp}}

\def\vH{\vec H}
\def\vwtH{\vec {\wt{H}}}
\newcommand{\Gp}[1]{\mathbf{G}^{(#1)}}
\newcommand{\gp}[1]{\mathbf{g}^{(#1)}}
\newcommand{\gb}[1]{\mathbf{g}^{[#1]}}
\newcommand{\wtHNp}[1]{{\wt{H}_N^{(#1)}}}

\newcommand{\HNp}[1]{{H_N^{(#1)}}}
\newcommand{\bgp}[1]{\bg^{(#1)}}

\newcommand{\Sum}{\mathrm{Sum}}
\def\depth{\text{depth}}
\def\type{\text{type}}
\def\Parent{\text{Parent}}



\newcommand{\norm}[1]{{\lt\|#1\rt\|}}
\newcommand{\tnorm}[1]{{\|#1\|}}
\newcommand{\bnorm}[1]{\big\|#1\big\|}

\newcommand{\He}{\mathrm{He}}
\newcommand{\tHe}{\widetilde \He}

\newcommand{\barh}{{\bar h}}

\newcommand{\oC}{\overline{C}}

\newcommand{\ucuL}{\underline{\cuL}}
\newcommand{\ocuL}{\overline{\cuL}}
\newcommand{\ocuK}{\overline{\cuK}}

\newcommand{\Bmin}{B_{\min}}

\newcommand{\uk}{\underline{k}}

\def\wtsH{{\wt{\mathscr{H}}}}
\def\wtcA{{\wt{\cA}}}
\def\st{{\mathrm{st}}}

\newcommand{\pderiv}[2]{{\fr{\partial #1}{\partial #2}}}

\newcommand{\kmax}{k_{\max}}
\newcommand{\Span}{\mathrm{span}}

\newcommand{\adj}{\mathrm{adj}}
\newcommand{\loc}{\mathrm{loc}}
\newcommand{\den}{\mathrm{den}}

\newcommand{\tp}{{\widetilde{p}}}
\newcommand{\wth}{{\widetilde{h}}}
\newcommand{\tPhi}{{\widetilde{\Phi}}}
\newcommand{\tDelta}{{\widetilde{\Delta}}}
\newcommand{\free}{\mathrm{free}}

\newcommand{\diff}[1]{{\mathrm{d}#1}}
\newcommand{\deriv}[1]{{\fr{\diff{}}{\diff{#1}}}}

\def\vchi{\vec{\chi}}
\def\vphi{\vec{\phi}}
\def\up{\underline{p}}
\def\ut{\underline{t}}
\def\uvphi{\underline{\vphi}}

\def\pplus{p_+}
\def\pminus{p_-}
\def\vphiplus{\vphi_+}
\def\vphiminus{\vphi_-}

\def\vQ{\vec Q}
\def\vR{\vec R}
\def\uvR{\underline{\vR}}

\def\oH{\overline H}
\def\uH{\underline H}
\def\oHess{\overline {\Hess}}
\def\Cov{\mathrm{Cov}}


% \def\psim{\stackrel{p}{\simeq}}
\def\psim{{\simeq}}
\def\km{D}
\def\ons{\mathbf{ons}}
\def\ONS{\mathbf{ONS}}
%\newcommand{\x}{\mathbf{x}}
%\newcommand{\z}{\mathbf{z}}
\def\X{{\boldsymbol{X}}}
\def\ZZ{{\boldsymbol{Z}}}
%\newcommand{\w}{\mathbf{w}}
\def\W{{\boldsymbol{W}}}
\def\A{{\boldsymbol{A}}}
\def\F{{\boldsymbol{F}}}
\def\hZZ{\widehat{\boldsymbol{Z}}}
\def\AMP{{\sf AMP}}
\def\LAMP{{\sf LAMP}}
\def\QQ{\mathcal{Q}}
\def\qq{{\boldsymbol{q}}}
\def\ww{{\boldsymbol{w}}}
\def\WW{{\boldsymbol{WW}}}
\def\VV{{\boldsymbol{V}}}
% \def\vv{{\boldsymbol{v}}}
\def\HH{{\boldsymbol{H}}}
\def\hh{{\boldsymbol{h}}}
\def\M{{\boldsymbol{m}}}


\def\bcG{{\boldsymbol \cal G}}
\def\Sym{{\rm Sym}}
\def\VF{{\mathcal J}}
\def\sp{{\sf p}}
\def\sP{{\sf P}}
\def\sE{{\sf E}}
\def\sVar{{\sf Var}}
% \def\hbm{{\hat{\boldsymbol m}}}
\def\tbm{{\tilde{\boldsymbol m}}}
\def\hm{\hat{m}}
\def\tH{\tilde{H}}
\def\bLambda{{\boldsymbol \Lambda}}
\def\spn{{\rm span}}

\def\hbA{\wh{\boldsymbol A}}
\def\tbA{\wt{\boldsymbol A}}
\def\hbz{\hat{\boldsymbol z}}

\def\bS{{\boldsymbol S}}
\def\bfone{{\boldsymbol 1}}

\def\err{{\sf err}}
\def\lambdaknot{\lambda_0}
\def\barR{\overline R}
\def\star{*}
\def\rhon{\rho^\infty}
\def\bfzero{\boldsymbol{0}}
\def\tcL{{\tilde {\mathcal L}}}
\def\tcT{{\tilde {\mathcal T}}}

\def\Lip{\mathrm{Lip}}
\def\Adm{\mathrm{Adm}}
\def\hAdm{\widehat{\mathrm{Adm}}}
\def\hM{\widehat M}
\def\oM{\overline M}

\def\CS{{\mathsf{CS}}}
\def\vDelta{{\vec\Delta}}

\newcommand{\tay}{{\mathrm{tay}}}


	\newtheorem{thm}{Theorem}[section]
	\newtheorem{prop}[thm]{Proposition}
	\newtheorem{cor}[thm]{Corollary}
	\newtheorem{lem}[thm]{Lemma}
	\newtheorem{defi}[thm]{Definition}
	\newtheorem{rmk}[thm]{Remark}
	\numberwithin{equation}{section}


\title{
Luce Permutation
}
\author{Mark Sellke \and Jiamin Wang}

\begin{document}

	\maketitle  

\section{Introduction}

    We first introduce the definition of the Luce model.
    Fix some positive real number $\theta_1,\cdots,\theta_n$ with total sum $w_n=\theta_1+\cdot+\theta_n$, consider the probability distribution on $\sigma\in \cS_n$:
    \begin{equation}\label{def:luce}
        p_n(\sigma_n) = \frac{\theta_{\sigma^{-1}(1)}}{w_n}\times\frac{\theta_{\sigma^{-1}(2)}}{w_n-\theta_{\sigma^{-1}(1)}}\times \frac{\theta_{\sigma^{-1}(3)}}{w_n-\theta_{\sigma^{-1}(1)}-\theta_{\sigma^{-1}(2)}} \cdots \times \frac{\theta_{\sigma^{-1}(n)}}{\theta_{\sigma^{-1}(n)}},
    \end{equation}
    we call this probability distribution on $\cS_n$ \emph{Luce distribution}, denoted by $\sigma\sim \Luce(\theta_1,\cdots,\theta_n)$.

\section{Fixed Point}

    The following permuton limit is proved in \cite{borga2025permutonlocallimitsluce}.
    \begin{thm}[Permuton limit]\label{thm:permuton_conv}
	Set $f_n(y):=\theta_{\lfloor yn \rfloor}$ for all $y\in[0,1]$ and assume that $f_n\to f$ almost everywhere for some positive, finite, measurable function $f$ on $[0,1]$.
	Let also $\mu$ be the law of the pair
	\begin{equation}\label{eq:limiting_rv}
            (X,Y):=\Big(U,F\big(E/f(U)\big)\Big),
        \end{equation}
	where $U$ is a uniform random variable on $[0,1]$, $E$ is an independent exponential random variable of parameter 1, and
	\begin{equation}\label{eq:defn-F}
		F(x):=1-\int_0^1\mathrm{e}^{-x f(y)}\,d y,\quad \text{for all $x\geq 0$.}
	\end{equation}
	If $\sigma_n\sim \Luce(\theta_1,\cdots,\theta_n)$, then the sequence $(\sigma_n)_n$ converges in probability in the permuton sense to the deterministic permuton $\mu$. 
    
    Moreover, the permuton $\mu$ is absolutely continuous w.r.t.\ the uniform permuton $\Leb_{[0,1]^2}$ and it has the following density:
	\begin{equation}\label{eq:density}
		\rho(x,y):=	\frac{f(x)\mathrm{e}^{-f(x)F^{-1}(y)}}{\int_0^1f(t)\mathrm{e}^{-f(t)F^{-1}(y)}\,d t}, \quad\text{for all $(x,y)\in(0,1)^2$},
	\end{equation}
    where $F$ is as in \eqref{eq:defn-F}.
    \end{thm}

    Roughly speaking, one would believe that $i$ is a fixed point with probability $n\rho(i/n,i/n)$. However, the number of fixed points is not a continuous function with respect to the permuton topology. In this section, we study the set of fixed points.

    \begin{thm}\label{thm:pointwise}
        With the setting of Theorem \ref{thm:permuton_conv}, for $\alpha,\beta \in (0,1)$ and two sequences $\{k_n\}_{n\geq 1}$ and $\{l_n\}_{n\geq 1}$ satisfying $k_n/n\to \alpha$, $f_n(k_n)\to f(\alpha)$ and $l_n/n\to \beta$ as $n\to \infty$, we have
        \begin{equation}
            \lim_{n\to \infty}n\mathbb P(\sigma^{-1}_n(k_n)=l_n)=\rho (\alpha,\beta).
        \end{equation}
    \end{thm}

    \begin{lem}\label{lem:rest-theta}
        With the setting in Theorem \ref{thm:permuton_conv}, for $\beta \in (0,1)$ and a sequence $\{l_n\}_{n\geq 1}$ satisfying $l_n/n\to \beta$ as $n\to \infty$, we have the following limit in probability:
        \begin{equation}
            \lim_{n\to \infty} \frac{1}{n}\sum_{i=1}^{n}\theta_i \mathbb I\{\sigma^{-1}_n(i)>l_n\} = \int_0^1f(t)e^{-f(t)F^{-1}(\beta)}dt.
        \end{equation}
    \end{lem}

    \begin{proof}
    Let $E_{n,1},\dots,E_{n,n}$ be independent exponential r.v.'s with $\mathbb{E}[E_{n,i}]=1/\theta_{i}$. By the exponential representation of the Luce model, we can express $\sigma_n(i)$ as
    \begin{equation}\label{eq:sigma-expression}
        \sigma^{-1}_n(i)=\sum_{j=1}^n\mathbf{1}\{E_{n,j}\le E_{n,i}\}.
    \end{equation}
    Note that for $j\neq i$,
    \[
    \mathbb{P}(E_{n,j}\le E_{n,i}\mid E_{n,i})=1-\mathrm{e}^{-\theta_{j}E_{n,i}}.
    \]
    Applying Hoeffding's inequality, we have
    \begin{equation}\label{eq:hoeffding}
        \mathbb P(|\sigma^{-1}_n(i)-Z_{n,i}|\ge t\mid E_{n,i})\le 2 e^{-t^2/2n},
    \end{equation}
    where $Z_{n,i}=\frac{1}{n}\sum_{j=1}^n(1-e^{\theta_jE_{n,i}})$. We take two continuous functions $g_1,g_2$ such that
    \[
    \mathbb I(x>\frac{l_n}{n}-n^{-\delta})\le g_1\le \mathbb I(x>\frac{l_n}{n}) \le g_2 \le \mathbb I(x>\frac{l_n}{n}+n^{-\delta})
    \]
    Now we fix a continuous function $g$ with $\|g\|_{L^{\infty}}\leq 1$, define
    \[
    T_n=\frac{1}{n}\sum_{i=1}^{n}\theta_i \mathbb I\{\sigma^{-1}_n(i)>l_n\}\quad \mbox{ and } \quad T_n'(g)=\frac{1}{n}\sum_{i=1}^{n}\theta_i g(Z_{n,i}).
    \]
    By \eqref{eq:hoeffding}, we have
    \begin{equation}\label{eq:sandwich}
        \mathbb P(T_n'(g_1)\leq T_n\leq T_n'(g_2))\to 1, \mbox{ as }n\to \infty.
    \end{equation}
    Moreover, by bounded convergence theorem we have
    \begin{align}
        \mathbb E T_n'(g)=\mathbb E \int_0^1f_n(x) g(Z_{n,\lfloor nx\rfloor})dx
        &=\int_0^\infty\int_0^1 f_n(x) g \lt(\int_0^1(1-\exp(-\frac{f_n(y)}{f_n(x)}z)\,dy\rt)dx\,\,e^{-z}dz\\
        &\to \int_0^\infty\int_0^1 f(x) g \lt(\int_0^1(1-\exp(-\frac{f(y)}{f(x)}z)\,dy\rt)dx\,\,e^{-z}dz 
        \\
        &=\int_0^1\int_0^1f(x)g(y)\rho(x,y)dxdy.
    \end{align}
    Meanwhile, we have
    \begin{align}
        \mathbb E T_n'(g)^2-(\mathbb ET_n'(g))^2= \frac{1}{n^2}\sum_{i=1}^n\theta_i^2 \lt[\mathbb E g(Z_{n,i})^2-(\mathbb Eg(Z_{n,i}))^2\rt]\leq \frac{1}{n^2}\sum_{i=1}^{n}\theta_i^2\leq \frac{\max_i\theta_i}{n}\to 0.
    \end{align}
    Therefore, we have $T_n'(g) \to \int_0^1\int_0^1f(x)g(y)\rho(x,y)dxdy$ in probability. Combined with \eqref{eq:sandwich}, we obtain that
    \begin{align}
        \lim_{n\to \infty} T_n =& \int_0^1\int_\beta^1f(x)\rho(x,y)dxdy 
        = \int_0^1 f(x) \int_\beta^1 \frac{f(x)e^{-f(x)F^{-1}(y)}}{\int_0^1 f(t)e^{-f(t)F^{-1}(y)}dt} dy\\
        \overset{y=F(u)}{=}& \int_0^1 f(x)\int_{F^{-1}(\beta)}^\infty \frac{f(x)e^{-f(x)u}}{\int_0^1 f(t)e^{-f(t)u}dt} \lt( \int_0^1 f(t)e^{-f(t)u}dt\rt) du
        = \int_0^1f(t)e^{-f(t)F^{-1}(\beta)}dt.
        \qedhere
    \end{align}
\end{proof}

\begin{lem}\label{lem:occurlate}
        With the setting in Theorem \ref{thm:permuton_conv}, for $\alpha,\beta \in (0,1)$ and two sequences $\{k_n\}_{n\geq 1}$ and $\{l_n\}_{n\geq 1}$ satisfying $k_n/n\to \alpha$, $f_n(k_n/n)\to f(\alpha)$, and $l_n/n\to \beta$ as $n\to \infty$, we have
        \begin{equation}
            \lim_{n\to \infty}\mathbb P (\sigma^{-1}_n(k_n)\ge l_n)=e^{-f(\alpha)F^{-1}(\beta)}
        \end{equation}
    \end{lem}
    \begin{proof}
        We keep using the notation $Z_{n,i}$ introduced in the proof of Lemma \ref{lem:rest-theta}.
        We have
        \begin{align}
            \mathbb P(\sigma^{-1}_n(k_n)\ge l_n)
            &\leq \mathbb P(Z_{n,i}> \frac{l_n}{n}-n^{-\delta})+2e^{n^{1-2\delta}/2}\\
            &=\mathbb P(\int_0^1(1-\exp(\frac{f_n(y)}{f_n(k_n/n)}E))dy>\frac{l_n}{n}+n^{-\delta})+2e^{n^{1-2\delta}/2}\\
            &\to \mathbb P(1-\int_0^1\exp (\frac{f(y)}{f(\alpha)}E)dy>\beta)=e^{-f(\alpha)F^{-1}(\beta)}.
        \end{align}
        Therefore, $\limsup_{n\to \infty }\mathbb P(\sigma^{-1}_n(k_n)\ge l_n)\le e^{-f(\alpha)F^{-1}(\beta)}$. Similarly, one could show that $\liminf_{n\to \infty }\mathbb P(\sigma^{-1}_n(k_n)\ge l_n)\ge e^{-f(\alpha)F^{-1}(\beta)}$, which completes the proof.
    \end{proof}

\begin{proof}[Proof of Theorem \ref{thm:pointwise}]
    Notice that 
    \begin{equation}
        n\mathbb P(\sigma^{-1}_n(k_n)=l_n)
        =\mathbb E\lt[\frac{\theta_{k_n}}{\frac{1}{n}\sum_{i=1}^n \theta_i \mathbb I\{\sigma^{-1}_n(i)>l_n\}}\cdot \mathbb I\{\sigma^{-1}_n(k_n)\ge l_n\}\rt].
    \end{equation}
    By \ref{lem:rest-theta}, we have
    \begin{equation}
        \frac{\theta_{k_n}}{\frac{1}{n}\sum_{i=1}^n \theta_i \mathbb I\{\sigma^{-1}_n(i)>l_n\}} \to \frac{f(\alpha)}{\int_0^1f(t)e^{-f(t)F^{-1}(\beta)}dt}, \mbox{ in probability as } n \to \infty.
    \end{equation}
    Combined with Lemma \ref{lem:occurlate}, we complete the proof.
\end{proof}

\subsection{Fixed Point Counting}
    Now we calculate the law of the number of fixed points. We first introduce a way to explore the fixed point set. For each $1\leq k\leq n$, let $\mathcal F_k$ be the sigma field generated by $(\sigma^{-1}(1),\cdots, \sigma^{-1}(k))$. Also define
    $p_k=\mathbb P(\sigma^{-1}(i)=i\mid \cF_{k-1})$. By definition we have
    \begin{equation}\label{def:pk}
        p_k=\frac{\theta_k\mathbb I\{\sigma^{-1}_n(k)\ge k\}}{\sum_{i=1}^n \theta_i \mathbb I\{\sigma^{-1}_n(i)\ge k\}} \in \cF_{k-1}.
    \end{equation}

    \begin{lem}\label{lem:sum-p-limit}
        With the setting of Theorem \ref{thm:permuton_conv} and $p_k$'s defined in \eqref{def:pk}, the following holds.
        \begin{enumerate}
            \item For any $\alpha \in [0,1)$,
            \begin{equation}\label{eq:p-firstmoment}
                \lim_{n\to \infty} \mathbb E  \sum_{k=1}^{\lfloor \alpha n \rfloor} p_k = \int_0^\alpha \rho(x,x) dx.
            \end{equation}
            \item For any $\alpha \in [0,1)$,
            \begin{equation}\label{eq:p-secondmoment}
                \lim_{n\to \infty} \mathbb E \lt( \sum_{k=1}^{\lfloor \alpha n \rfloor} p_k \rt) ^2= \lt( \int_0^\alpha \rho(x,x) dx \rt)^2.
            \end{equation}
            \item For any $\alpha \in [0,1)$,
            \begin{equation}\label{eq:max-p}
                \max_{1\leq k\leq \lfloor \alpha n\rfloor} p_k \to 0 \quad \mbox{ in probability as } n \to \infty.
            \end{equation}
        \end{enumerate}
        In particular, for any $\alpha\in [0,1)$ we have $\sum_{k=1}^{\lfloor \alpha n \rfloor} p_k\to \int_0^\alpha \rho(x,x) dx $ in probability.
    \end{lem}

    \begin{proof}
        We fixed an integer $N>0$ and $m_j=\lfloor j \alpha n/N \rfloor$ for $0\leq j\leq N$. 
        By Lemma \ref{lem:rest-theta}, we could take $n$ large enough such that the following holds with probability at least $1-1/N$:
        \begin{equation}
            \lt |\frac{1}{n}\sum_{i=1}^n \theta_i \mathbb I\{\sigma^{-1}_n(i)\ge m_j\}- \int_0^1f(t)e^{-f(t)F^{-1}(j\alpha/N)}dt\rt|\leq \frac{1}{N};
        \end{equation}
        For $1\leq j,j'\leq N$, $m_{j-1}<k\leq m_j$, $m_{j'-1}<k\le m_{j'}$, we have
        \begin{equation}\label{eq:Epk}
            \mathbb E p_k = \theta_k \cdot \mathbb E \frac{\mathbb I\{\sigma^{-1}_n(k)\ge k\}}{\sum_{i=1}^n \theta_i \mathbb I\{\sigma^{-1}_n(i)\ge k\}} \leq  \theta_k \cdot \mathbb E \frac{\mathbb I\{\sigma^{-1}_n(k)\ge m_{j-1}\}}{\sum_{i=1}^n \theta_i \mathbb I\{\sigma^{-1}_n(i)\ge m_{j}\}}\leq \frac{\theta_k}{n} \cdot\lt [ \frac{\mathbb P(\sigma_n^{-1}(k)\ge m_{j-1})}{\int_0^1f(t)e^{-f(t)F^{-1}(j\alpha/N)}dt} + \frac{C}{N} \rt];
        \end{equation}
        \begin{equation}\label{eq:Epk2}
            \mathbb E p_k ^2 \leq \frac{\theta_k^2}{n^2} \cdot\lt [ \frac{\mathbb P(\sigma_n^{-1}(k)\ge m_{j-1})}{(\int_0^1f(t)e^{-f(t)F^{-1}(j\alpha/N)}dt)^2} + \frac{C}{N} \rt].
        \end{equation}
        Similarly, when $k\neq k'$ we have
        \begin{equation}\label{eq:Epkpk'}
            \mathbb E p_k p_{k'} \leq \frac{\theta_k \theta_k'}{n^2} \cdot\lt [ \frac{\mathbb P(\sigma_n^{-1}(k)\ge m_{j-1};\sigma_n^{-1}(k')\ge m_{j'-1})}{\int_0^1f(t)e^{-f(t)F^{-1}(j\alpha/N)}dt\cdot \int_0^1f(t)e^{-f(t)F^{-1}(j'\alpha/N)}dt} + \frac{C}{N} \rt].
        \end{equation}
        Using the notation $Z_{n,k}$ introduced in Lemma \ref{lem:rest-theta}, we have by \eqref{eq:hoeffding} that
        \begin{align}
            \mathbb P(\sigma^{-1}_n(k)\ge m_{j-1})
            &\leq \mathbb P(Z_{n,k}> \frac{m_{j-1}}{n}-n^{-\delta})+2e^{n^{1-2\delta}/2}\\
            &=\mathbb P(\int_0^1(1-\exp(\frac{f_n(y)}{\theta_k}E))dy>\frac{m_{j-1}}{n}+n^{-\delta})+2e^{n^{1-2\delta}/2}\\
            &\le e^{-\theta_k F_n^{-1}((j-1)\alpha/N-n^{-\delta})}+2e^{n^{1-2\delta}/2}.\label{eq:1notoccur}
        \end{align}
        where $F_n(x)=1-\int_0^1 (1- e^{-xf_n(y)})dy$. Similarly, we have
        \begin{equation}\label{eq:2notoccur}
            \mathbb P(\sigma^{-1}_n(k)\ge m_{j-1};\sigma^{-1}_n(k')\ge m_{j'-1})\leq  e^{-\theta_k F_n^{-1}((j-1)\alpha/N-n^{-\delta})} e^{-\theta_{k'} F_n^{-1}((j'-1)\alpha/N-n^{-\delta})}+4e^{n^{1-2\delta}/2}.
        \end{equation}
        
        We first prove \eqref{eq:p-firstmoment}. Combining \eqref{eq:Epk} and \eqref{eq:1notoccur}, we get
        \begin{align}
            \sum_{k=1}^{\lfloor \alpha n \rfloor} \mathbb Ep_k 
            & \leq \sum_{j=1}^N \sum_{k=m_{j-1}+1}^{m_j} \frac{f_n(k/n)e^{-f_n(k/n) F_n^{-1}((j-1)\alpha/N-n^{-\delta})}}{n\int_0^1f(t)e^{-f(t)F^{-1}(j\alpha/N)}dt} + \frac{C}{N}\\
            &=\sum_{j=1}^N  \int_{m_{j-1}/n}^{m_j\alpha/N} \frac{f_n(x)e^{-f_n(x) F_n^{-1}((j-1)\alpha/N-n^{-\delta})}}{\int_0^1f(t)e^{-f(t)F^{-1}(j\alpha/N)}dt} \, dx + \frac{C}{N}.
        \end{align}
        where we used the fact that $\sum_{k=1}^{\lfloor \alpha n \rfloor}\theta_k\le n$.
        Since $f_n(x)\to f(x)$, we see that $F_n(x)\to F(x)$ and $F_n^{-1}(x)\to F^{-1}(x)$ almost everywhere. Thus we obtain by bounded convergence theorem that
        \begin{equation}\label{eq:1stmoment-upper}
            \limsup_{n\to \infty} \sum_{k=1}^{\lfloor \alpha n \rfloor} \mathbb Ep_k 
            \leq \sum_{j=1}^N  \int_{(j-1)\alpha/N}^{j\alpha/N} \frac{f(x)e^{-f(x) F^{-1}((j-1)\alpha/N)}}{\int_0^1f(t)e^{-f(t)F^{-1}(j\alpha/N)}dt} \, dx + \frac{1}{N^2} = \int_0^\alpha \frac{f(x)e^{-f(x)F^{-1}(\lfloor Nx\rfloor/N)}}{\int_0^1f(t)e^{-f(t)F^{-1}(\lceil Nx\rceil/N)}dt} \,dx +\frac{1}{N^2}.
        \end{equation}
        Similarly, we get that
        \begin{equation}\label{eq:1stmoment-lower}
            \liminf_{n\to \infty} \sum_{k=1}^{\lfloor \alpha n \rfloor} \mathbb Ep_k 
            \geq \int_0^\alpha \frac{f(x)e^{-f(x)F^{-1}(\lceil Nx\rceil/N)}}{\int_0^1f(t)e^{-f(t)F^{-1}(\lfloor Nx\rfloor/N)}dt} \,dx -\frac{1}{N^2}.
        \end{equation}
        Note that 
        \begin{equation}
            \frac{f(x)e^{-f(x)F^{-1}(\lceil Nx\rceil/N)}}{\int_0^1f(t)e^{-f(t)F^{-1}(\lfloor Nx\rfloor/N)}dt} \leq \rho (x,x) \leq \frac{f(x)e^{-f(x)F^{-1}(\lfloor Nx\rfloor/N)}}{\int_0^1f(t)e^{-f(t)F^{-1}(\lceil Nx\rceil/N)}dt},
        \end{equation}
        \begin{equation}
            \lim_{N\to \infty} \frac{f(x)e^{-f(x)F^{-1}(\lfloor Nx\rfloor/N)}}{\int_0^1f(t)e^{-f(t)F^{-1}(\lceil Nx\rceil/N)}dt} \,dx = \lim_{N\to \infty}\frac{f(x)e^{-f(x)F^{-1}(\lceil Nx\rceil/N)}}{\int_0^1f(t)e^{-f(t)F^{-1}(\lfloor Nx\rfloor/N)}dt} =\rho(x,x),
        \end{equation}
        using monotone convergence theorem, we conclude by taking $N\to \infty $ on both sides of \eqref{eq:1stmoment-upper} and \eqref{eq:1stmoment-lower}.

        Next, we prove \eqref{eq:p-secondmoment}. Combining \eqref{eq:Epk2}, \eqref{eq:Epkpk'} and \eqref{eq:2notoccur}, we get
        \begin{align}
            &\mathbb E\lt(\sum_{k=1}^{\lfloor \alpha n \rfloor} p_k \rt)^2 = \sum_{k\neq k'} \mathbb Ep_kp_{k'}+ \sum_{k=1}^n \mathbb Ep_k^2 \\
            & \leq \sum_{j,j'=1}^N \sum_{k=m_{j-1}+1}^{m_j} \sum_{k'=m_{j'-1}+1}^{m_{j'}} \frac{f_n(k/n)e^{-f_n(k/n) F_n^{-1}((j-1)\alpha/N-n^{-\delta})}}{n\int_0^1f(t)e^{-f(t)F^{-1}(j\alpha/N)}dt} \cdot \frac{f_n(k'/n)e^{-f_n(k'/n) F_n^{-1}((j'-1)\alpha/N-n^{-\delta})}}{n\int_0^1f(t)e^{-f(t)F^{-1}(j'\alpha/N)}dt}
            \\
            &+\sum_{j=1}^N \frac{\theta_k^2e^{-f_n(k/n) F_n^{-1}((j-1)\alpha/N-n^{-\delta})}}{n^2(\int_0^1f(t)e^{-f(t)F^{-1}(j\alpha/N)}dt)^2} + \frac{1}{N^2} \\
            &=\lt(\sum_{j=1}^N  \int_{m_{j-1}/n}^{m_j\alpha/N} \frac{f_n(x)e^{-f_n(x) F_n^{-1}((j-1)\alpha/N-n^{-\delta})}}{\int_0^1f(t)e^{-f(t)F^{-1}(j\alpha/N)}dt} \, dx \rt)^2 +\sum_{k=1}^N \frac{\theta_k^2e^{-f_n(k/n) F_n^{-1}((j-1)\alpha/N-n^{-\delta})}}{n^2(\int_0^1f(t)e^{-f(t)F^{-1}(j\alpha/N)}dt)^2} + \frac{1}{N^2}.
        \end{align}
        Notice that $\alpha<1$, we have
        \begin{equation}
            \int_0^1f(t)e^{-f(t)F^{-1}(j\alpha/N)}dt \geq \int_0^1f(t)e^{-f(t)F^{-1}(\alpha)}dt.
        \end{equation}
        Therefore, we have
        \begin{equation}
            \sum_{j=1}^N \frac{\theta_k^2e^{-f_n(k/n) F_n^{-1}((j-1)\alpha/N-n^{-\delta})}}{n^2(\int_0^1f(t)e^{-f(t)F^{-1}(j\alpha/N)}dt)^2} \leq \frac{1}{\int_0^1f(t)e^{-f(t)F^{-1}(\alpha)}dt} \cdot \sum_{j=1}^N \frac{\theta_k^2}{n^2}\leq \frac{\max_k \theta_k/n}{\int_0^1f(t)e^{-f(t)F^{-1}(\alpha)}dt} 
        \end{equation}
        Taking $n\to \infty$, we get that
        \begin{equation}
            \limsup_{n\to \infty} \mathbb E\lt(\sum_{k=1}^{\lfloor \alpha n \rfloor} p_k \rt)^2  \leq \lt( \int_0^\alpha \frac{f(x)e^{-f(x)F^{-1}(\lfloor Nx\rfloor/N)}}{\int_0^1f(t)e^{-f(t)F^{-1}(\lceil Nx\rceil/N)}dt} \,dx \rt)^2 +\frac{1}{N^2}.
        \end{equation}
        On the other hand, by Cauchy-Schwarz inequality, we have
        \begin{equation}\label{eq:2ndmoment-upper}
            \mathbb E\lt(\sum_{k=1}^{\lfloor \alpha n \rfloor} p_k \rt)^2  \geq \lt(\mathbb E\sum_{k=1}^{\lfloor \alpha n \rfloor} p_k \rt)^2 .
        \end{equation}
        By \eqref{eq:1stmoment-lower}, we get that
        \begin{equation}\label{eq:2ndmoment-lower}
            \limsup_{n\to \infty} \mathbb E\lt(\sum_{k=1}^{\lfloor \alpha n \rfloor} p_k \rt)^2  \leq \lt( \int_0^\alpha \frac{f(x)e^{-f(x)F^{-1}(\lceil Nx\rceil/N)}}{\int_0^1f(t)e^{-f(t)F^{-1}(\lfloor Nx\rfloor/N)}dt} \,dx  -\frac{1}{N^2}\rt)^2.
        \end{equation}
        Similar to the first moment, we conclude by taking $N\to \infty $ on both sides of \eqref{eq:2ndmoment-upper} and \eqref{eq:2ndmoment-lower}.
    \end{proof}

    Now we introduce the following Poisson limit theorem for martingales.

    \begin{lem}
        Consider a probability space $(\Omega,\mathcal F, \mathbb P)$ with a series of filtration $\{\cF_{n,i}\}_{n\geq 1, 1\leq i\leq m(n)}$. Suppose that $p_{n,i}\in \cF_{n,i-1}$ are random variables taking values in $[0,1]$, satisfying:
        \begin{enumerate}
            \item $p_{n,1}+\cdots + p_{n,m(n)} \to \lambda$ as $n \to \infty$, in probability;
            \item $\max_{1\leq i\leq m(n)}p_{n,i}\to 0$ as $n \to \infty$, in probability.
        \end{enumerate}
        For each $n\geq 1, 1\leq i\leq m(n)$, let $U_{n,i}$ be an independent random variable with uniform distribution on $[0,1]$. Let  $X_n=\sum_{i=1}^{m(n)} \mathbb I\{U_{n,i}<p_{n,i}\}$, we have $X_n\to Poi(\lambda)$ as $n\to \infty$ in distribution.
    \end{lem}

    \subsection{Fixed Point Locations}
    In this subsection we study the probability that $k_1<\cdots< k_l\le \alpha n$ are the only fixed points. We call this event $A_{k_1,\cdots k_l}$. Let $\mu_n$ be the law of the fixed point set.
    Moreover, let $\nu_n$ be the measure on subset of $[\alpha n]$ such that for each $k\in[n]$, whether it belongs to the set is determined by an independent Bernoulli random variable with parameter
    \[q_k=\frac{\theta_{k}e^{-\theta_{k}F_n^{-1}(k/n)}}{n\int_0^1f(t)e^{-f(t)F^{-1}(k/n)}}.\]
    Then we have
    \begin{equation}\label{eq:q-limit}
        \sum_{1\leq k\leq \alpha n}q_k \to \int_0^\alpha \rho(x,x)\,dx \mbox{ and }\max_{1\leq k\leq \alpha n}q_k \to 0.
    \end{equation}

    \begin{lem}
        We have 
        \begin{equation}
            \lim_{n\to \infty} d_{\TV} (\mu_n,\nu_n) =0.
        \end{equation}
    \end{lem}

    \begin{proof}
        Fix $N,L,\eps$. 
        First, since $f_n(x)\to f(x)$ almost everywhere, there exists $C(\eps)$, such that by taking $n$ sufficiently large, 
        \begin{equation}
            \sum_{k\in A(K)} \theta_k \leq \eps n \mbox{ where } A(K)=\{k:\theta_k\geq K\}.
        \end{equation}
        We take $N\gg K(\eps)$. By \eqref{eq:q-limit} we have
        \begin{equation}\label{eq:nu-nound}
            \nu_{n}(\{k_1,\cdots,k_l\})=q_{k_1}\cdots q_{k_l}\prod_{t\neq k_j} (1-q_{t})\leq \prod_{1\leq i\leq l} \frac{\theta_{k_i}e^{-\theta_{k_i}F_n^{-1}(k_i/n)}}{n\int_0^1f(t)e^{-f(t)F^{-1}(k_i/n)}} \exp\lt(-\int_0^\alpha \rho (x,x) \,dx \rt) (1+\frac{1}{N^{1/2}})
        \end{equation}       
        Notice that by the definition of $p_k$'s as in \eqref{def:pk}, we have
        \begin{equation}\label{eq:onlykifixed}
            \mu_n(\{k_1,\cdots k_l\})=\mathbb E \lt[ p_{k_1}\cdots p_{k_l} \prod_{t\neq k_i, 1\le t\le \alpha n} (1-p_t) \rt].
        \end{equation}
        Since by Lemma \ref{lem:sum-p-limit} we have in probability, as $n \to \infty$,
        \begin{equation}
            \sum_{1\le t\le \alpha n} p_t \to \int_0^\alpha \rho (x,x) \,dx .
        \end{equation}
        Combined with \eqref{eq:max-p}, we could take $n$ large enough such that the following holds with probability at least $1-1/N$: for each $l \leq L$ and $1\leq k_1<\cdots<k_l\leq \alpha n$,
        \begin{equation}\label{eq:other-not-fixed}
            \lt | \prod_{t\neq k_i, 1\le t\le \alpha n} (1-p_t) - \exp\lt(-\int_0^\alpha \rho (x,x) \,dx \rt) \rt | \leq \frac{1}{N} .
        \end{equation}
        We could then deal with $\mathbb  E p_{k_1}\cdots p_{k_l}$ as above. Let $m_j=\lfloor j\alpha n/N\rfloor$By Lemma \ref{lem:rest-theta}, we could take $n$ large enough such that the following holds with probability at least $1-1/N$: for each $1\leq j\leq N$
        \begin{equation}\label{eq:restestimate}
            \lt |\frac{1}{n}\sum_{i=1}^n \theta_i \mathbb I\{\sigma^{-1}_n(i)\ge m_j\}- \int_0^1f(t)e^{-f(t)F^{-1}(j\alpha/N)}dt\rt|\leq \frac{1}{N};
        \end{equation}
        Therefore, let $\Omega'\subset \Omega$ such that $\mathbb P(\Omega')\geq 1-2/N$, and \eqref{eq:other-not-fixed} and \eqref{eq:restestimate} hold. Let $\mu_n'$ be the law of the fixed point set restricted on $\Omega'$. Also let $m_{j_i-1}<k_i\leq m_{j_i}$, and we have
        \begin{equation}\label{eq:kifixed}
            \mathbb E[p_{k_1}\cdots p_{k_l};\Omega']\geq \mathbb P(\sigma^{-1}(k_i)\geq m_{j_i}, \forall 1\leq j\leq k;\Omega')\cdot \frac{1}{n^l}\lt[\prod_{1\leq i\leq l}\frac{\theta_{k_i}}{\int_0^1f(t)e^{-f(t)F^{-1}((j_i-1)\alpha/N)}dt}\rt] \cdot (1-\frac{L}{N})
        \end{equation}
        Similar to \eqref{eq:2notoccur}, we have
        \begin{equation}
            \mathbb P(\sigma^{-1}_n(k_i)\ge m_{j_i}\forall 1\leq i\leq l;\Omega')\geq \prod_{1\leq i\leq l}e^{-\theta_{k_i}F_n^{-1}(k_i/n)} -Le^{-n^{1-2\delta}/2} -\frac{2}{N}.
        \end{equation}
        Thus, when $k_1,\cdots,k_l \notin A(K)$, since $N\gg K(\eps)$, we have
        \begin{equation}\label{eq:lnotoccur}
            \mathbb P(\sigma^{-1}_n(k_i)\ge m_{j_i}\forall 1\leq i\leq l;\Omega')\geq \prod_{1\leq i\leq l}e^{-\theta_{k_i}F_n^{-1}(k_i/n)}(1-\frac{1}{N^{1/2}}).
        \end{equation}
        We deduce from \eqref{eq:onlykifixed}, \eqref{eq:other-not-fixed}, \eqref{eq:kifixed} and \eqref{eq:lnotoccur} that, by taking $n$ sufficiently large depending on $N,L$, the following holds: for each $1\leq l\leq L$ and $1\leq k_1<\cdots<k_{l}\leq \alpha n$ with $k_1,\cdots,k_l \notin A(K)$,
        \begin{equation}\label{eq:mu-bound}
            \mu_n(\{k_1,\cdots,k_l\}) \geq \prod_{1\leq i\leq l} \frac{\theta_{k_i}e^{-\theta_{k_i}F_n^{-1}(k_i/n)}}{n\int_0^1f(t)e^{-f(t)F^{-1}(k_i/n-1/N)}} \exp\lt(-\int_0^\alpha \rho (x,x) \,dx \rt) (1-\frac{L}{N^{1/2}}).
        \end{equation}
        Taking $N$ sufficiently large w.r.t. $\eps$, we have
        \begin{equation}\label{eq:Nlarge}
            \lt| \int_0^1f(t)e^{-f(t)F^{-1}(u-1/N)}-\int_0^1f(t)e^{-f(t)F^{-1}(u)} \rt|\leq \eps \mbox{ and } \frac{L}{N^{1/2}}\leq \eps.
        \end{equation}
        Combining \eqref{eq:nu-nound}, \eqref{eq:mu-bound} and \eqref{eq:Nlarge}, we have
        \begin{align*}
            \mu_n(\{k_1,\cdots,k_l\})
            &\geq \prod_{1\leq i\leq l} \frac{\theta_{k_i}e^{-\theta_{k_i}F_n^{-1}(k_i/n)}}{n\int_0^1f(t)e^{-f(t)F^{-1}(k_i/n)}} \exp\lt(-\int_0^\alpha \rho (x,x) \,dx \rt) (1-2L\eps)\\
            &\geq (1-3L\eps) \nu_n(\{k_1,\cdots,k_l\}).
        \end{align*}
        While for $\{k_1,\cdots,k_l\}\cap A(K)\neq \emptyset$, we have
        \begin{align*}
            \sum_{\{k_1,\cdots,k_l\}\cap A(K)\neq \emptyset} \nu_n(\{k_1,\cdots,k_l\}) 
            &\leq \sum_{\{k_1,\cdots,k_l\}\cap A(K)\neq \emptyset} \frac{\theta_{k_1}\cdots \theta_{k_l}}{n^l}\cdot\lt( \frac{1}{\int_{0}^1 f(t)e^{-f(t)} F^{-1}(\alpha)}\rt)^l
            \\
            &\leq L\cdot \frac{\sum_{k\in A(K)}\theta_k}{n}\cdot\lt( \frac{1}{\int_{0}^1 f(t)e^{-f(t)} F^{-1}(\alpha)}\rt)^l
            \\
            &\leq
            L\eps\cdot\lt( \frac{1}{\int_{0}^1 f(t)e^{-f(t)} F^{-1}(\alpha)}\rt)^l.
        \end{align*}        
        Thus we have
        \begin{align*}
            &\sum_{l=1}^L\sum_{k_1<\cdots<k_l}(\nu_n(\{k_1,\cdots,k_l\})-\mu_n(\{k_1,\cdots,k_l\}))_+ \\
            &\leq 3L\eps\sum_{l=1}^L\sum_{k_1<\cdots<k_l}\nu_n(\{k_1,\cdots,k_l\})+\sum_{l=1}^LL\eps\cdot\lt( \frac{1}{\int_{0}^1 f(t)e^{-f(t)} F^{-1}(\alpha)}\rt)^l\leq C(L)\eps.
        \end{align*}
        Therefore, taking $\eps < 1/LC(L)$, we have
        \begin{align}
            d_{\TV}(\mu_n,\nu_n)&\leq \sum_{l=1}^L\sum_{k_1<\cdots<k_l}(\nu_n(\{k_1,\cdots,k_l\})-\mu_n(\{k_1,\cdots,k_l\}))_+ + \sum_{l>L}\sum_{k_1<\cdots<k_l}(\nu_n(\{k_1,\cdots,k_l\})
            \\
            &\leq \frac{1}{L} + \sum_{l>L}\sum_{k_1<\cdots<k_l}\nu_n(\{k_1,\cdots,k_l\}.
        \end{align}
        Notice that $\nu_n\to Poi(\int_0^\alpha \rho(x,x)\,dx)$ in distribution, we have $\sum_{l>L}\sum_{k_1<\cdots<k_l}\nu_n(\{k_1,\cdots,k_l\})=o_L(1)$. We therefore conclude by taking $L$ to infinity. 
    \end{proof}



% \section{Total Variation Poisson Limit}


% Let's assume for concreteness that $\theta_i\in [1/C,C]$ for a uniform constant $C$.
% Then we claim the set of fixed points is within $o(1)$ total variation distance of a product measure $\mu_{prod}$ on $\{1,2,\dots,n\}$ which has probability $\rho(i/n,i/n)/n$ to include $i$.
% (And the technical condition $\theta_i\in [1/C,C]$ should rule out fixed points in the last $o(n)$ steps, with high probability.)
% Proof outline:

% \begin{itemize}
%     \item 
%     Truncate to ``typical'' permutations that are within $o(1)$ of the limiting permuton.
%     Call this subprobability measure $\mu_{typ}$.
%     \item
%     Fix $k=O(1)$ and $\eps>0$.
%     Given any $k_1<i_2\dots<k_l\leq (1-\eps)n$, let $E$ be the event that these are all fixed points, and $\hat E$ the event that these are the only fixed points.
%     We want to show that 
%     \[
%     \mu_{typ}(\hat E)/\mu_{prod}(\hat E)
%     =
%     1\pm o(1).
%     \]
%     \item 
%     To show this, we use the same filtration as above.
%     In each step we need the next point not to be fixed, except at times $k_1<i_2\dots<k_l$, and we need $(k_1,...,k_l)$ to not appear when they are not supposed to.
%     \item 
%     At times $k_1,...,k_l$, as long as the permuton limit holds, the probability to be a fixed point will be 
%     \[
%     (1\pm o(1))\cdot\rho(i/n,i/n)/n.
%     \]
%     \item 
%     At other times say between $k_i$ and $i_{j+1}$, the probability to be a fixed point will be something like
%     \[
%     \frac{1-h(i/n,i/n)-h(i_{j+1}/n,i/n)-\dots-h(k_l/n,i/n)\pm o(1)}{n}.
%     \]
%     (Here $h(x)$ is some nice functionin terms of $f$ and/or $\rho$ that I am not sure of the exact formula for. I am just trying to say that we just have to avoid the elements $(i,i_{j+1},\dots,k_l)$ and I think there should be a simple formula for this conditional probability.) 
%     \begin{itemize}
%     \item
%     More precisely, the $h(i/n,i/n)$ term will not always appear, depending on whether $i$ has been revealed yet, and if it does appear it will be larger.
%     But the ``local average'' value of this term will be something like $h(i/n,i/n)$ as long as the permuton limit holds, which is all we need.
%     Namely on $i/n \in [\alpha,\alpha+\delta]$ the number of $\rho(i/n,i/n)$ terms is a simple continuous function of 
%     \[
%     \frac{|\{j<\alpha n:\theta(j)<\alpha n\}|}{n}
%     =
%     \frac{|\{j<(\alpha+\delta) n:\theta(j)<(\alpha+\delta) n\}|}{n} \pm O(\delta)
%     \]
%     which is itself a continuous function on permuton space.
%     \end{itemize}
%     \item 
%     There is something slightly annoying to argue because in the likelihood ratios above, we don't want to deal with the truncation of $\mu_{typ}$.
%     I think one way is to say that if we truncate at most $n^{-10}$ probability to get $\mu_{typ}$, then with probability $1-n^{-5}$, the conditional probability under the original $\mu$ (given the past) to end up being typical is $\geq 1-n^{-5}$ (by Doob's inequality). 
%     So on this event, all values above are ``uniformly approximately correct'' and the truncation shouldn't cause any more issues.
% \end{itemize}


\mscomment{GPT output below, with some edits.}

\section{Late Fixed Point Counting}

Let's assume that:
\begin{enumerate}
    \item The empirical distribution $\frac{1}{n}\sum_{i=1}^n \delta_{\theta_i}$ is tight in $(0,\infty)$. Informally, ``nearly all values are $\Theta(1)$''.
    \item 
    We have $\theta_m\leq C<\infty$ for all $m\leq \eps n$.
    \item 
    We have $\theta_{n-m}\geq \gamma>0$ for all  $m\leq \eps n$.
\end{enumerate}

We will show the expected number of fixed points near the end is $o(1)$.
The plan is:

\begin{enumerate}[label=(\alph*)]
    \item If $\theta_i\geq (\log n)^2$ then $\theta_i$ appears early.
    \label{it:1}
    \item 
    The last $n^{o(1)}$ steps have $o(1)$ points on average, even if the last $n^{o(1)}$ are ordered in an adversarial way.
    \item 
    For $n^{o(1)}\ll m\ll n$, the number and total weight of the active exponential variables concentrate.
    \item 
    Then we can just compute everything almost exactly, and end up with a bound by a telescoping sum of the increments of a monotone function.
\end{enumerate}

\begin{lem}[Very large rates are selected early, hence can be truncated]
\label{lem:huge-early}
Let $(E_i)_{i=1}^n$ be independent exponentials with rates $(\theta_i)_{i=1}^n$,
and let $\sigma^{-1}(i)$ be the rank of $E_i$ among $(E_1,\dots,E_n)$ (so $\sigma^{-1}(i)$
is the position of item $i$ in the Luce ordering).

Assume a \emph{tightness} condition: there exists a finite $K$ such that for all
sufficiently large $n$,
\[
\bigl|\{i\in[n]: \theta_i \le K\}\bigr| \ge \tfrac34\,n .
\]
Let $L_n := (\log n)^2$. Then
\[
\mathbb P\Bigl(\exists i\in[n]\ \text{with}\ \theta_i \ge L_n\ \text{and}\ \sigma^{-1}(i) > \tfrac n2\Bigr)
\le e^{-c n} + n e^{-t_0 (\log n)^2} = o(1),
\]
where one may take $t_0 := \frac1K\log\frac43$ and $c=c(K)>0$.

In particular, with probability $1-o(1)$, every index with $\theta_i \ge L_n$ appears
among the first $n/2$ positions of the permutation.
\end{lem}

\begin{proof}
Define the alive-count process
\[
N(t) := \sum_{i=1}^n \mathbf 1\{E_i>t\},\qquad t\ge 0.
\]
Let $B:=\{i:\theta_i\le K\}$, so $|B|\ge 3n/4$.
Set $t_0 := \frac1K\log\frac43$, so that $e^{-K t_0}=\frac34$.

For each $i\in B$, we have $\mathbb P(E_i>t_0)=e^{-\theta_i t_0}\ge e^{-K t_0}=\frac34$.
Hence
\[
\mathbb E\bigl[N(t_0)\bigr] \;\ge\; \mathbb E\Bigl[\sum_{i\in B}\mathbf 1\{E_i>t_0\}\Bigr]
\;\ge\; |B|\cdot \frac34 \;\ge\; \frac{9}{16}n.
\]
By a Chernoff bound for a sum of independent Bernoulli variables,
\[
\mathbb P\bigl(N(t_0) < n/2\bigr) \le e^{-c n}
\]
for some $c=c(K)>0$.

Let $T_{n/2}$ denote the (random) time when the alive-count first drops to $n/2$
(equivalently, $T_{n/2}$ is the $n/2$-th order statistic among $E_1,\dots,E_n$).
On the event $\{N(t_0)\ge n/2\}$ we must have $T_{n/2}\ge t_0$.

Now fix any $i$ with $\theta_i\ge L_n$. Then
\[
\mathbb P(E_i>t_0)=e^{-\theta_i t_0}\le e^{-L_n t_0}.
\]
By a union bound over at most $n$ indices,
\[
\mathbb P\bigl(\exists i:\theta_i\ge L_n \ \text{and}\ E_i>t_0\bigr)\le n e^{-L_n t_0}.
\]
Finally, on $\{N(t_0)\ge n/2\}$ we have $T_{n/2}\ge t_0$, so
\[
\{\theta_i\ge L_n \ \text{and}\ \sigma^{-1}(i)>n/2\}
\ \subseteq\ 
\{\theta_i\ge L_n \ \text{and}\ E_i>T_{n/2}\}
\ \subseteq\
\{\theta_i\ge L_n \ \text{and}\ E_i>t_0\}.
\]
Combining with the two bounds above yields the claim.
\end{proof}

\subsection{A subpolynomial tail is negligible}
\label{subsec:subpoly}

Let $m=m(n)$ with $m=n^{o(1)}$ and define $\alpha_n\coloneqq 1-m/n$.
Write
\[
N_n(1)=N_n(\alpha_n)+T_n,
\qquad
T_n\coloneqq \sum_{k=\lfloor \alpha_n n\rfloor+1}^n \mathbf 1\{\sigma_n^{-1}(k)=k\}.
\]
Suppose all the weights are at least $\gamma$ on $[(1-\varepsilon)n,n]$),
the expected number of $k\geq n-m$ such that $\sigma_n(k)\geq n-m$ is $o(1)$.
Indeed with very high probability at least $\sqrt{n}$ of the middle $N/3$ exponential variables will be $\geq \frac{\eps\log n}{C'}$ by tightness of the empirical measure \ref{it:1}. 
Meanwhile the average number of the last $m$ exponential variables which are $\geq \frac{\eps\log n}{C}$ is at most 
\[
me^{\Omega(\log n)}=mn^{-\Omega(1)}\ll 1.
\]



     
     \subsection{Tail fixed points via quantile zoom and a telescoping bound}

     Given Subsection~\ref{subsec:subpoly} and Lemma~\ref{lem:huge-early}, it remains to count fixed points which are in the last $o(n)$ but not the last $n^{o(1)}$, and we can ignore any $\theta$ values larger than $(\log n)^2$.
    This will let us show concentration, at any time $T$, for the number of remaining exponentials (that are still bigger than $T$) as well as the total weight of all the remaining ones. 
    (Focusing on times where the typical number of remaining points is $\gg n^{o(1)}$ is key for concentration, and ignoring small $\theta$ values is needed for concentration of the sum.)

\textbf{Setup and normalization.}
Throughout this subsection we use the exponential representation of the Luce model:
let $(E_i)_{i=1}^n$ be independent exponentials with $\mathbb E[E_i]=1/\theta_i$.
Then $\sigma^{-1}(i)$ is the rank of $E_i$ among $\{E_j\}_{j=1}^n$, i.e.
\begin{equation}\label{eq:exp-repr}
\sigma^{-1}(i)=\sum_{j=1}^n \mathbf 1\{E_j\le E_i\}.
\end{equation}
Since scaling all weights by a constant does not change the Luce law, we may (and do)
assume the convenient normalization
\begin{equation}\label{eq:normalize}
\sum_{i=1}^n \theta_i = n.
\end{equation}

For $1\le k\le n$, let $\mathcal F_k:=\sigma(\sigma^{-1}(1),\dots,\sigma^{-1}(k))$ and define
\begin{equation}\label{eq:pk-def}
p_k := \mathbb P(\sigma^{-1}(k)=k \mid \mathcal F_{k-1})
     = \frac{\theta_k \mathbf 1\{\sigma^{-1}(k)\ge k\}}{\sum_{i=1}^n \theta_i \mathbf 1\{\sigma^{-1}(i)\ge k\}}
     \in \mathcal F_{k-1}.
\end{equation}
Then
\begin{equation}\label{eq:tail-exp-sum-pk}
\mathbb E\big|\{k:\sigma^{-1}(k)=k,\ k\in[n-r_n+1,n]\}\big|
= \sum_{k=n-r_n+1}^n \mathbb E[p_k].
\end{equation}

\textbf{Continuous-time viewpoint and deterministic quantile times.}
Define the (random) survivor count and survivor total weight at time $t\ge0$:
\begin{equation}\label{eq:NtWt}
N(t):=\sum_{i=1}^n \mathbf 1\{E_i>t\}, \qquad
W(t):=\sum_{i=1}^n \theta_i \mathbf 1\{E_i>t\}.
\end{equation}
Also define the corresponding deterministic mean curves
\begin{equation}\label{eq:HnDn}
H_n(t):=\frac1n\sum_{i=1}^n e^{-\theta_i t}, \qquad
D_n(t):=\frac1n\sum_{i=1}^n \theta_i e^{-\theta_i t},
\end{equation}
so that $\mathbb E[N(t)]=nH_n(t)$ and $\mathbb E[W(t)]=nD_n(t)$.
Note that $H_n$ is strictly decreasing with $H_n(0)=1$ and $H_n(t)\downarrow0$.

For each integer $m\in\{1,\dots,n\}$, define the deterministic \emph{quantile time} $t_m$ by
\begin{equation}\label{eq:tm-def}
H_n(t_m)=\frac{m}{n}.
\end{equation}
Thus $t_m$ is the deterministic time at which the \emph{expected} number of survivors equals $m$.

Let $T_m$ be the (random) time at which the number of survivors first drops to $m$:
\begin{equation}\label{eq:Tm-def}
T_m:=\inf\{t\ge0:\ N(t)\le m\}.
\end{equation}
At the discrete step $k=n-m+1$, the set of remaining labels is exactly $\{i:\ E_i>T_m\}$,
and hence from \eqref{eq:pk-def} we may rewrite
\begin{equation}\label{eq:pk-Tm}
p_{n-m+1}=\frac{\theta_{n-m+1}\mathbf 1\{E_{n-m+1}>T_m\}}{W(T_m)}.
\end{equation}


\textbf{A high-probability ``quantile window'' event.}
Fix a small tail proportion $\varepsilon\in(0,1/2)$ and set $M:=\lfloor \varepsilon n\rfloor$.
We will analyze $p_{n-m+1}$ for 
\[
(\log n)^{10}\le m\le M.
\]
Let $\eta\in(0,1/10)$ be a small accuracy parameter to be chosen later.

For each $m$ define the window size $r_m:=\lceil m^{2/3}\rceil$.
Consider the event $\mathcal G_n(\varepsilon,\eta)$ that simultaneously for all $m\in\{1,\dots,M\}$,
\begin{equation}\label{eq:G-event}
\big|N(t_m)-m\big|\le r_m
\qquad\text{and}\qquad
\big|W^{(\le L_n)}(t_m)-nD_n^{(\le L_n)}(t_m)\big|\le \eta\,nD_n^{(\le L_n)}(t_m),
\end{equation}
where $D_n^{(\le L_n)}(t):=\frac1n\sum_{\theta_i\le L_n}\theta_i e^{-\theta_i t}$.
Since $N(t_m)$ is a sum of independent Bernoullis with mean $m$ and $W^{(\le L_n)}(t_m)$ is a sum of independent
bounded variables in $[0,L_n]$, a Chernoff--Bernstein bound and a union bound over $m\le M$ give:
for every fixed $\varepsilon,\eta>0$,
\begin{equation}\label{eq:G-prob}
\mathbb P(\mathcal G_n(\varepsilon,\eta))\ge 1-o(1)\qquad(n\to\infty).
\end{equation}

On $\mathcal G_n(\varepsilon,\eta)$, monotonicity of $N(\cdot)$ implies the random time $T_m$ is trapped between
deterministic quantile times:
\begin{equation}\label{eq:Tm-bracket}
t_{m+r_m}\ \le\ T_m\ \le\ t_{m-r_m}.
\end{equation}

\textbf{Pointwise upper bound on $\mathbb E[p_{n-m+1}]$.}
Fix $m\le M$ and set $k:=n-m+1$.
Using \eqref{eq:pk-Tm}, \eqref{eq:Tm-bracket}, and monotonicity of $W^{(\le L_n)}(\cdot)$, on $\mathcal G_n$ we have
\begin{align}
p_k
&=\frac{\theta_k\mathbf 1\{E_k>T_m\}}{W(T_m)}
\ \le\ \frac{\theta_k\mathbf 1\{E_k>t_{m+r_m}\}}{W^{(\le L_n)}(T_m)}
\ \le\ \frac{\theta_k\mathbf 1\{E_k>t_{m+r_m}\}}{W^{(\le L_n)}(t_{m-r_m})}. \label{eq:pk-basic-bound}
\end{align}
On $\mathcal G_n$, the denominator satisfies
$W^{(\le L_n)}(t_{m-r_m})\ge (1-\eta)\,nD_n^{(\le L_n)}(t_{m-r_m})$, hence
\begin{equation}\label{eq:pk-good}
p_k\,\mathbf 1_{\mathcal G_n}
\ \le\ \frac{\theta_k}{(1-\eta)n}\cdot \frac{\mathbf 1\{E_k>t_{m+r_m}\}}{D_n^{(\le L_n)}(t_{m-r_m})}.
\end{equation}
Taking expectations and using $\mathbb P(E_k>t)=e^{-\theta_k t}$ yields
\begin{equation}\label{eq:Ekpk-preasymp}
\mathbb E[p_k]
\ \le\ \frac{\theta_k e^{-\theta_k t_{m+r_m}}}{(1-\eta)n\,D_n^{(\le L_n)}(t_{m-r_m})}
\ +\ \mathbb P(\mathcal G_n^c).
\end{equation}

Now assume the tail lower bound $\theta_k\ge\gamma$ holds for all $k\in[(1-\varepsilon)n,n]$.
Since $m\le \varepsilon n$ implies $k=n-m+1\in[(1-\varepsilon)n,n]$, we have $\theta_k\ge\gamma$.
Moreover, by taking $\varepsilon$ sufficiently small (depending on $\gamma$ and the tightness constant used below),
we may ensure $t_{m+r_m}\ge 1/\gamma$ for all $m\le M$ and all large $n$; on this range
$\theta\mapsto \theta e^{-\theta t}$ is decreasing for $\theta\ge\gamma$, hence
\begin{equation}\label{eq:theta-lower-replace}
\theta_k e^{-\theta_k t_{m+r_m}} \le \gamma e^{-\gamma t_{m+r_m}}.
\end{equation}
Plugging \eqref{eq:theta-lower-replace} into \eqref{eq:Ekpk-preasymp} gives the clean estimate
\begin{equation}\label{eq:Ekpk-gamma}
\mathbb E[p_{n-m+1}]
\ \le\ \frac{\gamma e^{-\gamma t_{m+r_m}}}{(1-\eta)n\,D_n^{(\le L_n)}(t_{m-r_m})}
\ +\ o(1),
\qquad \text{uniformly for }1\le m\le M.
\end{equation}

\textbf{From pointwise bounds to a telescoping sum.}

\mscomment{Here one might need to argue some regularity properties. I think it should be fine though because $H_n$ and $D_n$ should only vary by $1\pm o(1)$ factors on $t_m\pm r_m$.}

Summing \eqref{eq:Ekpk-gamma} over $m=1,\dots,M$ and using \eqref{eq:tail-exp-sum-pk} gives
\begin{equation}\label{eq:sumEkpk-mid}
\mathbb E\big|\{k:\sigma^{-1}(k)=k,\ k\in[(1-\varepsilon)n,n]\}\big|
\ \le\ \frac{1}{1-\eta}\sum_{m=1}^{M}\frac{\gamma e^{-\gamma t_{m+r_m}}}{n\,D_n^{(\le L_n)}(t_{m-r_m})}
\ +\ o(1).
\end{equation}

To connect to the ``ODE/quantile'' time-change, introduce the continuous parameter $s\in(0,1]$ and define
$t(s):=H_n^{-1}(s)$ so that $t(m/n)=t_m$.
Define $G(s):=e^{-\gamma t(s)}$.
A direct differentiation using $H_n'(t)=-D_n(t)$ shows
\begin{equation}\label{eq:Gprime}
G'(s)=\frac{\gamma e^{-\gamma t(s)}}{D_n(t(s))}.
\end{equation}
Thus, heuristically, the main term in \eqref{eq:sumEkpk-mid} is a Riemann sum for
$\int_0^\varepsilon G'(s)\,ds = G(\varepsilon)-G(0)=e^{-\gamma t(\varepsilon)}=e^{-\gamma t_M}$.
The discrete indices $m\pm r_m$ simply fatten the mesh; since $r_m=o(m)$ uniformly for $m\to\infty$,
this does not change the limiting integral once $m$ is not extremely small.
(For the very last $m\le n^{o(1)}$ one can simply keep the crude bound \eqref{eq:Ekpk-gamma}
and note it already forces a negligible total contribution; see the remark below.)

In particular, combining \eqref{eq:Gprime} with a standard Riemann-sum comparison (on $[n^{-1/3},\varepsilon]$)
yields the deterministic ``telescoping'' bound:
\begin{equation}\label{eq:telescoping-sum}
\sum_{m=1}^{M}\frac{\gamma e^{-\gamma t_m}}{n\,D_n(t_m)}
\ \le\ e^{-\gamma t_M}\ +\ o(1).
\end{equation}
Using that $D_n^{(\le L_n)}(t)=D_n(t)+o(1)$ uniformly on the relevant $t$-range (by tightness/truncation),
and that $t_{m\pm r_m}=t_m+o(1)$ for $m\le M$ once $m\to\infty$, we may replace the shifted/truncated
version in \eqref{eq:sumEkpk-mid} by \eqref{eq:telescoping-sum} at the cost of $o(1)$.
Therefore, for fixed $\varepsilon>0$ and $\eta>0$,
\begin{equation}\label{eq:tailbound-to-exp}
\mathbb E\big|\{k:\sigma^{-1}(k)=k,\ k\in[(1-\varepsilon)n,n]\}\big|
\ \le\ \frac{1}{1-\eta}\,e^{-\gamma t_M}\ +\ o(1).
\end{equation}

\textbf{Converting $e^{-\gamma t_M}$ into a power of $\varepsilon$ using tightness.}
Recall we assume the weight profile is tight, so for every $\delta>0$ there exists $K<\infty$
such that for all large $n$,
\begin{equation}\label{eq:tightness-K}
\big|\{i:\theta_i\le K\}\big|\ge (1-\delta)n.
\end{equation}
Then for all $t\ge0$,
\[
H_n(t)=\frac1n\sum_{i=1}^n e^{-\theta_i t}\ \ge\ \frac1n\sum_{\theta_i\le K} e^{-\theta_i t}
\ \ge\ (1-\delta)e^{-Kt}.
\]
Plugging $t=t_M$ and $H_n(t_M)=M/n\approx\varepsilon$ gives
\begin{equation}\label{eq:tM-lower}
t_M \ \ge\ \frac1K\log\frac{1-\delta}{\varepsilon},
\qquad\text{hence}\qquad
e^{-\gamma t_M}\ \le\ \Big(\frac{\varepsilon}{1-\delta}\Big)^{\gamma/K}.
\end{equation}
Combining \eqref{eq:tailbound-to-exp} and \eqref{eq:tM-lower} yields:
for any fixed $\delta,\eta>0$, all sufficiently large $n$ satisfy
\begin{equation}\label{eq:tail-final}
\mathbb E\big|\{k:\sigma^{-1}(k)=k,\ k\in[(1-\varepsilon)n,n]\}\big|
\ \le\ \frac{1}{1-\eta}\Big(\frac{\varepsilon}{1-\delta}\Big)^{\gamma/K}\ +\ o(1).
\end{equation}
Since $\delta,\eta$ are arbitrary, this gives the claimed $\varepsilon^c+\text{(small)}$ bound,
with $c=\gamma/K$.


\mscomment{asked for more details on the ODE approximation and got:}


% =====================  Late tail fixed points: last o(n)  =====================

\subsection{No fixed points in the last $o(n)$ positions}

In this subsection we prove that, under a lower bound on the weights in the last $r_n=o(n)$ indices,
there are no fixed points there with high probability.

\begin{thm}[Last $o(n)$ indices contain no fixed points w.h.p.]\label{thm:tail_on}
Let $r_n=o(n)$ and let $\gamma>0$.
Assume that
\begin{equation}\label{eq:tail-lower}
\theta_k \ge \gamma\qquad\text{for all }k\in\{n-r_n+1,\dots,n\}.
\end{equation}
Then
\[
\mathbb P\Big(\exists\,k\in\{n-r_n+1,\dots,n\}\text{ such that }\sigma_n^{-1}(k)=k\Big)\to 0,
\qquad n\to\infty.
\]
Equivalently,
\[
\mathbb E\Big|\{k:\ \sigma_n^{-1}(k)=k,\ k\in\{n-r_n+1,\dots,n\}\}\Big|\to 0.
\]
\end{thm}

\begin{proof}
\textbf{Step 0: normalization and exponential race.}
By scale invariance of the Luce distribution, we may assume
\begin{equation}\label{eq:normalize-sumtheta}
\sum_{i=1}^n \theta_i = n.
\end{equation}
Let $E_1,\dots,E_n$ be independent exponentials with rates $\theta_1,\dots,\theta_n$.
Then $\sigma_n^{-1}(i)$ is the rank of $E_i$ among $E_1,\dots,E_n$.

\textbf{Step 1: tail fixed points are controlled by $\sum \mathbb E[p_k]$.}
Recall the exploration filtration $\mathcal F_k=\sigma(\sigma_n^{-1}(1),\dots,\sigma_n^{-1}(k))$ and
\[
p_k:=\mathbb P(\sigma_n^{-1}(k)=k\mid \mathcal F_{k-1})
=\frac{\theta_k \mathbf 1\{\sigma_n^{-1}(k)\ge k\}}{\sum_{i=1}^n \theta_i \mathbf 1\{\sigma_n^{-1}(i)\ge k\}}
\in \mathcal F_{k-1}.
\]
Then
\begin{equation}\label{eq:tail-exp-sum}
\mathbb E\Big|\{k:\sigma_n^{-1}(k)=k,\ k\in\{n-r_n+1,\dots,n\}\}\Big|
=\sum_{k=n-r_n+1}^{n}\mathbb E[p_k].
\end{equation}

\textbf{Step 2: deterministic quantile times and the ODE.}
Define for $t\ge0$ the (random) survivor count and survivor weight:
\[
N(t):=\sum_{i=1}^n \mathbf 1\{E_i>t\},\qquad
W(t):=\sum_{i=1}^n \theta_i\,\mathbf 1\{E_i>t\}.
\]
Define the deterministic mean curves
\begin{equation}\label{eq:HnDn-def}
H_n(t):=\frac1n\sum_{i=1}^n e^{-\theta_i t},\qquad
D_n(t):=\frac1n\sum_{i=1}^n \theta_i e^{-\theta_i t}.
\end{equation}
Then $\mathbb E[N(t)]=nH_n(t)$ and $\mathbb E[W(t)]=nD_n(t)$, and $H_n'(t)=-D_n(t)$.

For each integer $m\in\{1,\dots,n\}$ define the deterministic quantile time $t_m$ by
\begin{equation}\label{eq:tm-def}
H_n(t_m)=\frac{m}{n}.
\end{equation}
Also define the random ``$m$ survivors time''
\[
T_m:=\inf\{t\ge0:\ N(t)\le m\}.
\]
At the discrete step $k=n-m+1$ there are exactly $m$ labels remaining, namely $\{i:E_i>T_m\}$, and
\begin{equation}\label{eq:pk-Tm}
p_{n-m+1}
=\frac{\theta_{n-m+1}\,\mathbf 1\{E_{n-m+1}>T_m\}}{W(T_m)}.
\end{equation}

\begin{lem}[ODE for the quantile map]\label{lem:ode}
Let $t(\cdot)$ be the inverse of $H_n$, i.e.\ $t(s):=H_n^{-1}(s)$ for $s\in(0,1]$.
Then $t$ is $C^1$ and satisfies
\begin{equation}\label{eq:ode-t}
t'(s)=-\frac{1}{D_n(t(s))}.
\end{equation}
Consequently, for $G(s):=e^{-\gamma t(s)}$ we have
\begin{equation}\label{eq:ode-G}
G'(s)=\frac{\gamma e^{-\gamma t(s)}}{D_n(t(s))}.
\end{equation}
\end{lem}

\begin{proof}
Differentiate the identity $H_n(t(s))=s$ to get $H_n'(t(s))\,t'(s)=1$.
Since $H_n'(t)=-D_n(t)$, this yields \eqref{eq:ode-t}. Then \eqref{eq:ode-G} follows by the chain rule.
\end{proof}

\textbf{Step 3: a deterministic telescoping bound.}
The next inequality is the deterministic heart of the tail estimate.

\begin{lem}[Telescoping bound]\label{lem:telescoping}
Let $(t_m)$ be defined by \eqref{eq:tm-def}. Then for every $M\in\{1,\dots,n\}$,
\begin{equation}\label{eq:telescoping}
\sum_{m=1}^{M}\frac{\gamma e^{-\gamma t_m}}{nD_n(t_m)} \ \le\ e^{-\gamma t_M}.
\end{equation}
\end{lem}

\begin{proof}
Using $H_n'(t)=-D_n(t)$ and $H_n(t_{m-1})-H_n(t_m)=\frac{1}{n}$, we have
\[
\frac{1}{n}
= \int_{t_m}^{t_{m-1}} D_n(u)\,du
\le D_n(t_m)\,(t_{m-1}-t_m)
\qquad\text{since }D_n\text{ is decreasing}.
\]
Therefore $\frac{1}{nD_n(t_m)}\le t_{m-1}-t_m$.
Multiplying by $\gamma e^{-\gamma t_m}$ and summing over $m\le M$ gives
\[
\sum_{m=1}^{M}\frac{\gamma e^{-\gamma t_m}}{nD_n(t_m)}
\le
\sum_{m=1}^{M}\gamma e^{-\gamma t_m}(t_{m-1}-t_m)
\le
\int_{t_M}^{\infty}\gamma e^{-\gamma t}\,dt
=e^{-\gamma t_M},
\]
where the last inequality uses that $t\mapsto e^{-\gamma t}$ is decreasing.
\end{proof}

\textbf{Step 4: concentration and bracketing of $T_m$.}
Fix a large constant $A>10$ and define a (deterministic) deviation level
\[
\Delta_m := \lceil A\sqrt{m}\log n\rceil,\qquad m\ge1.
\]
For $m\ge (\log n)^4$ we have $\Delta_m\le m/2$ for all large $n$.

\begin{lem}[Bracketing $T_m$ by deterministic quantiles]\label{lem:bracket-Tm}
There exists $A_0<\infty$ such that, for the choice $\Delta_m=\lceil A_0\sqrt m\log n\rceil$,
\[
\mathbb P\big(\forall m\in[(\log n)^4,M]:\ t_{m+\Delta_m}\le T_m\le t_{m-\Delta_m}\big)\ \ge\ 1-o(1).
\]
\end{lem}

\begin{proof}
Fix $m$ and write $t^-:=t_{m+\Delta_m}$ and $t^+:=t_{m-\Delta_m}$.
By monotonicity of $N(\cdot)$, we have the equivalences
\[
\{T_m\ge t\}\iff \{N(t)\ge m\},\qquad \{T_m\le t\}\iff \{N(t)\le m\}.
\]
Note that $\mathbb E[N(t^-)]=nH_n(t^-)=m+\Delta_m$ and $\mathbb E[N(t^+)]=m-\Delta_m$.

Since $N(t)$ is a sum of independent Bernoullis, a Chernoff bound yields
\[
\mathbb P\big(N(t^-)<m\big)\le \exp\Big(-c\,\frac{\Delta_m^2}{m+\Delta_m}\Big),
\qquad
\mathbb P\big(N(t^+)>m\big)\le \exp\Big(-c\,\frac{\Delta_m^2}{m}\Big),
\]
for a universal $c>0$. With $\Delta_m\asymp \sqrt m\log n$, both probabilities are $\le n^{-c' \log n}$.
A union bound over $m\le M$ (with $M\le n$) gives the claim.
\end{proof}

\textbf{Step 5: lower-bounding $W(T_m)$ by $nD_n(t_{m-\Delta_m})$ on a good event.}
Fix $L_n:=(\log n)^2$ and let $W^{(\le L_n)}(t):=\sum_{\theta_i\le L_n}\theta_i\mathbf 1\{E_i>t\}$ and
$D_n^{(\le L_n)}(t):=\frac1n\sum_{\theta_i\le L_n}\theta_i e^{-\theta_i t}$.
(If you already have a lemma stating that indices with $\theta_i>L_n$ are selected before time $n/2$
w.h.p., you may replace $W$ by $W^{(\le L_n)}$ for tail steps $k>n/2$.)

\begin{lem}[Weight concentration at the bracket time]\label{lem:Wt-conc}
Fix $\eta\in(0,1/2)$. With probability $1-o(1)$, simultaneously for all
$m\in[(\log n)^4,M]$,
\begin{equation}\label{eq:Wt-lb}
W^{(\le L_n)}(t_{m-\Delta_m}) \ \ge\ (1-\eta)\,nD_n^{(\le L_n)}(t_{m-\Delta_m}).
\end{equation}
\end{lem}

\begin{proof}
For each fixed $t$, $W^{(\le L_n)}(t)$ is a sum of independent random variables taking values in $[0,L_n]$.
A multiplicative Chernoff--Bernstein bound gives
\[
\mathbb P\Big(W^{(\le L_n)}(t)\le (1-\eta)\mathbb E W^{(\le L_n)}(t)\Big)
\le \exp\Big(-c\,\frac{\eta^2\,\mathbb E W^{(\le L_n)}(t)}{L_n}\Big),
\]
for a universal $c>0$.
Apply this at $t=t_{m-\Delta_m}$, for which $\mathbb E W^{(\le L_n)}(t)=nD_n^{(\le L_n)}(t)$.
Since $m\ge(\log n)^4$ and $\Delta_m\asymp \sqrt m\log n$, one has $H_n(t_{m-\Delta_m})=(m-\Delta_m)/n\ge (\log n)^4/(2n)$,
and by \eqref{eq:normalize-sumtheta} and $\theta_i\le L_n$ on the truncation set,
\[
nD_n^{(\le L_n)}(t)\ge \sum_{\theta_i\le L_n}\theta_i e^{-\theta_i t}
\ge e^{-L_n t}\sum_{\theta_i\le L_n}\theta_i
\ge 0.
\]
(If desired, one may further restrict to $m\ge (\log n)^{10}$ so that $nD_n^{(\le L_n)}(t_{m-\Delta_m})\gg L_n\log n$ deterministically,
making the above tail probability super-polynomial in $n$ uniformly over $m$; the remaining smaller $m$ contribute $o(1)$ to the tail count
by the last step of the proof below.)
A union bound over $m\le M$ completes the argument.
\end{proof}

\textbf{Step 6: pointwise bound on $\mathbb E[p_{n-m+1}]$.}
Fix $m\in[(\log n)^4,M]$ and write $k=n-m+1$. On the event from Lemmas~\ref{lem:bracket-Tm} and \ref{lem:Wt-conc},
we have $T_m\le t_{m-\Delta_m}$ and hence (since $W^{(\le L_n)}$ is decreasing in $t$)
\[
W(T_m)\ \ge\ W^{(\le L_n)}(T_m)\ \ge\ W^{(\le L_n)}(t_{m-\Delta_m})
\ \ge\ (1-\eta)\,nD_n^{(\le L_n)}(t_{m-\Delta_m}),
\]
and also $T_m\ge t_{m+\Delta_m}$ implies $\mathbf 1\{E_k>T_m\}\le \mathbf 1\{E_k>t_{m+\Delta_m}\}$.
Thus, using \eqref{eq:pk-Tm},
\begin{equation}\label{eq:pk-bound-shifted}
p_k \ \le\ \frac{\theta_k}{(1-\eta)n}\cdot
\frac{\mathbf 1\{E_k>t_{m+\Delta_m}\}}{D_n^{(\le L_n)}(t_{m-\Delta_m})}
\qquad\text{on the good event.}
\end{equation}
Taking expectations and using $\mathbb P(E_k>t)=e^{-\theta_k t}$ gives
\begin{equation}\label{eq:Epk-bound-shifted}
\mathbb E[p_{n-m+1}]
\ \le\ \frac{\theta_{n-m+1}\,e^{-\theta_{n-m+1}t_{m+\Delta_m}}}{(1-\eta)n\,D_n^{(\le L_n)}(t_{m-\Delta_m})}
\ +\ o(1/M),
\end{equation}
uniformly for $m\in[(\log n)^4,M]$.

\textbf{Step 7: smoothness in $m$ (removing the $\pm\Delta_m$ shifts).}
We now justify replacing $t_{m\pm\Delta_m}$ and $D_n^{(\le L_n)}(t_{m-\Delta_m})$ by $t_m$ and $D_n^{(\le L_n)}(t_m)$.
This is the only place where we use an ``ODE/smoothness'' input.

\begin{lem}[Quantile smoothness]\label{lem:smoothness}
Assume we are on the truncation event $\{\theta_i\le L_n \text{ for all } i\in\{n-M+1,\dots,n\}\}$ (this holds w.h.p.\ if $M=o(n)$ under tightness).
Fix $m\in[(\log n)^{10},M]$. Then
\begin{equation}\label{eq:tm-smooth}
t_m-t_{m+\Delta_m}\ \le\ \frac{\Delta_m}{nD_n(t_m)}.
\end{equation}
Moreover, $D_n$ is log-Lipschitz in time:
\begin{equation}\label{eq:Dn-loglip}
\bigl|\log D_n(t)-\log D_n(s)\bigr|\ \le\ L_n\,|t-s|\qquad\forall s,t\ge0,
\end{equation}
and hence, whenever $L_n\,(t_m-t_{m+\Delta_m})=o(1)$,
\begin{equation}\label{eq:Dn-smooth}
D_n(t_{m+\Delta_m})=(1+o(1))D_n(t_m)
\quad\text{and}\quad
e^{-\theta_k t_{m+\Delta_m}}=(1+o(1))e^{-\theta_k t_m}
\ \ \text{for all }\theta_k\le L_n.
\end{equation}
The same statements hold with $(m+\Delta_m)$ replaced by $(m-\Delta_m)$.
\end{lem}

\begin{proof}
For \eqref{eq:tm-smooth}, note that
\[
\frac{\Delta_m}{n}
=H_n(t_{m+\Delta_m})-H_n(t_m)
=\int_{t_{m+\Delta_m}}^{t_m} D_n(u)\,du
\ge D_n(t_m)\,(t_m-t_{m+\Delta_m}),
\]
since $u\le t_m$ on the integration range and $D_n$ is decreasing.

For \eqref{eq:Dn-loglip}, differentiate:
\[
D_n'(t)= -\frac1n\sum_{i=1}^n \theta_i^2 e^{-\theta_i t},
\]
and on the truncation set $\theta_i\le L_n$ we have $\theta_i^2\le L_n\theta_i$, hence
$|D_n'(t)|\le L_n D_n(t)$, which implies \eqref{eq:Dn-loglip} by integrating $|(\log D_n)'|\le L_n$.

Finally, if $L_n|t-s|=o(1)$, then \eqref{eq:Dn-loglip} implies $D_n(t)/D_n(s)=1+o(1)$.
The exponential smoothness follows similarly from $e^{-\theta(t-s)}=1+o(1)$ when $\theta\le L_n$ and $L_n|t-s|=o(1)$.
\end{proof}

Combining Lemma~\ref{lem:smoothness} with \eqref{eq:Epk-bound-shifted} shows that, for $m\in[(\log n)^{10},M]$,
\begin{equation}\label{eq:Epk-main}
\mathbb E[p_{n-m+1}]
\ \le\ (1+o(1))\cdot \frac{\theta_{n-m+1}e^{-\theta_{n-m+1}t_m}}{nD_n(t_m)}\ +\ o(1/M).
\end{equation}

\textbf{Step 8: impose the tail lower bound and sum using the telescoping lemma.}
Assume \eqref{eq:tail-lower}. For $m\le M$, $k=n-m+1$ lies in the last $M$ indices, hence $\theta_k\ge\gamma$.
Also, since $M=o(n)$, we have $t_M\to\infty$ and in particular $t_m\ge 1/\gamma$ for all $m\le M$ and all large $n$.
On $[1/\gamma,\infty)$ the function $a\mapsto a e^{-a t}$ is decreasing for $a\ge \gamma$, so
\[
\theta_k e^{-\theta_k t_m}\ \le\ \gamma e^{-\gamma t_m}.
\]
Therefore, summing \eqref{eq:Epk-main} over $m\in[(\log n)^{10},M]$ and using Lemma~\ref{lem:telescoping} gives
\[
\sum_{m=(\log n)^{10}}^{M}\mathbb E[p_{n-m+1}]
\ \le\ (1+o(1))\sum_{m=1}^{M}\frac{\gamma e^{-\gamma t_m}}{nD_n(t_m)} + o(1)
\ \le\ (1+o(1))e^{-\gamma t_M}+o(1).
\]
The remaining $m<(\log n)^{10}$ contribute $o(1)$ to \eqref{eq:tail-exp-sum} because $M=o(n)$ and each term is at most $1$,
and (more sharply) because $\mathbb P(E_k>t_M)\le e^{-\gamma t_M}$ and there are only $(\log n)^{10}$ such indices.

Thus, using \eqref{eq:tail-exp-sum},
\begin{equation}\label{eq:tail-E-bound}
\mathbb E\Big|\{k:\sigma_n^{-1}(k)=k,\ k\in\{n-M+1,\dots,n\}\}\Big|
\ \le\ C\,e^{-\gamma t_M}+o(1),
\end{equation}
for some absolute $C<\infty$.

\textbf{Step 9: lower bound $t_M$ and conclude.}
By \eqref{eq:normalize-sumtheta}, at least $n/2$ indices satisfy $\theta_i\le 2$, hence for all $t\ge0$,
\[
H_n(t)=\frac1n\sum_{i=1}^n e^{-\theta_i t}\ \ge\ \frac12 e^{-2t}.
\]
Since $H_n(t_M)=M/n$, we obtain $M/n \ge \frac12 e^{-2t_M}$ and therefore
\[
e^{-\gamma t_M}\ \le\ (2M/n)^{\gamma/2}.
\]
If $M=r_n=o(n)$, then $(2M/n)^{\gamma/2}\to0$.
Plugging into \eqref{eq:tail-E-bound} yields
\[
\mathbb E\Big|\{k:\sigma_n^{-1}(k)=k,\ k\in\{n-r_n+1,\dots,n\}\}\Big|\to0,
\]
and Markov's inequality gives the desired high-probability statement.
\end{proof}
% ============================================================================


     
     \footnotesize
\bibliographystyle{alphaabbr}
\bibliography{lucebib}
\end{document}



